\begin{proof}[Proof of \cref{obj:thm:adaptive-minimax}]
\emph{1. Uniform local approximation.}
We first establish the estimator bound uniformly over the prescribed adaptation interval. Set
\[
m:=\lceil\beta\rceil-1,\qquad
s:=\beta-m,\qquad
\mathcal A_{\mathrm{poly}}
:=\{\alpha\in\mathbb N_0^d:|\alpha|\le m\},
\qquad
p_{\mathrm{coef}}:=|\mathcal A_{\mathrm{poly}}|.
\]
Then \(0<s\le1\) and \(p_{\mathrm{coef}}>0\). By the identification and uniqueness conclusions of \cref{obj:lem:causal-identification}, every Hölder response witnessing membership of \(P\) in \(\mathcal P_\gamma\) agrees with \(\mu_{1,P}\) throughout \(\mathcal X\).

Here is the local approximation bound used below, including at cube boundaries. Define the affine map and transported response by
\[
T_{\mathrm{aff}}(x):=2x-\mathbf1
\quad(x\in\mathbb R^d),
\qquad
g_{\mathrm{aff}}(z)
:=\mu_{1,P}\!\left(\frac{z+\mathbf1}{2}\right)
\quad(z\in[-1,1]^d).
\]
Repeated coordinate differentiation of the affine composition gives, for every multiindex \(\alpha\) with \(|\alpha|\le m\),
\[
D^\alpha g_{\mathrm{aff}}(z)
=
2^{-|\alpha|}D^\alpha\mu_{1,P}\!\left(\frac{z+\mathbf1}{2}\right),
\qquad z\in[-1,1]^d.
\]
The chain rule proves this identity in the interior, and continuity of the derivative traces extends it to the boundary. Thus \cref{obj:ass:holder-derivative-bound} gives
\[
\sup_{z\in[-1,1]^d}|D^\alpha g_{\mathrm{aff}}(z)|
\le 2^{-|\alpha|}L\le L.
\]
For \(|\alpha|=m\) and \(z,z'\in[-1,1]^d\), the identity and \cref{obj:ass:holder-modulus} give
\[
\begin{aligned}
|D^\alpha g_{\mathrm{aff}}(z)-D^\alpha g_{\mathrm{aff}}(z')|
&=2^{-m}\left|
D^\alpha\mu_{1,P}\!\left(\frac{z+\mathbf1}{2}\right)
-D^\alpha\mu_{1,P}\!\left(\frac{z'+\mathbf1}{2}\right)
\right|
\\
&\le 2^{-m}L
\left\|\frac{z-z'}2\right\|_\infty^s
\\
&\le L\|z-z'\|_\infty^s.
\end{aligned}
\]
Affine composition also preserves the continuous derivative traces through order \(m\). Consequently, defining
\[
K_{\mathrm{aff}}:=1,
\]
we obtain intrinsic Hölder radius \(K_{\mathrm{aff}}L\) for \(g_{\mathrm{aff}}\) on \([-1,1]^d\). \Cref{obj:lem:global-holder-extension}, applied with derivative order \(m\), exponent \(s\), and radius \(K_{\mathrm{aff}}L\), supplies a constant \(K_{\mathrm{ext}}>0\) and a function \(g_{\mathrm{ext}}\in C^m(\mathbb R^d)\) satisfying
\[
g_{\mathrm{ext}}(z)=g_{\mathrm{aff}}(z)
\quad(z\in[-1,1]^d),
\]
\[
\max_{0\le k_{\mathrm{der}}\le m}
\sup_{z\in\mathbb R^d}
\|\mathrm D^{k_{\mathrm{der}}}g_{\mathrm{ext}}(z)\|_{\mathrm{op}}
\le K_{\mathrm{ext}}K_{\mathrm{aff}}L,
\]
and
\[
\|\mathrm D^m g_{\mathrm{ext}}(z)
-\mathrm D^m g_{\mathrm{ext}}(z')\|_{\mathrm{op}}
\le
K_{\mathrm{ext}}K_{\mathrm{aff}}L
\|z-z'\|_\infty^s
\quad(z,z'\in\mathbb R^d).
\]
Here \(\mathrm D^{k_{\mathrm{der}}}\) denotes the Fréchet derivative of the indicated order, with order zero interpreted as the function itself.

For a partition cube \(Q\) of mesh \(0<h\le1\), define
\[
\vartheta_{Q,\alpha}
:=
\frac{(2h)^{|\alpha|}}{\alpha!}
D^\alpha g_{\mathrm{ext}}(T_{\mathrm{aff}}(a_Q)),
\qquad \alpha\in\mathcal A_{\mathrm{poly}},
\]
and
\[
\Pi_Q(x)
:=
U\!\left(\frac{x-a_Q}{h}\right)^\top\vartheta_Q,
\qquad x\in Q.
\]
We derive the remainder bound along the expansion segment. Since \(\beta>1\), we have \(m\ge1\). For \(x\in Q\), define
\[
z_{Q,0}:=T_{\mathrm{aff}}(a_Q),
\qquad
v_{Q,x}:=2(x-a_Q),
\qquad
\varphi_{Q,x}(t_{\mathrm{seg}})
:=g_{\mathrm{ext}}(z_{Q,0}+t_{\mathrm{seg}}v_{Q,x})
\quad(t_{\mathrm{seg}}\in\mathbb R).
\]
The chain rule gives
\[
\varphi_{Q,x}^{(k_{\mathrm{der}})}(t_{\mathrm{seg}})
=
\mathrm D^{k_{\mathrm{der}}}g_{\mathrm{ext}}
(z_{Q,0}+t_{\mathrm{seg}}v_{Q,x})
[\underbrace{v_{Q,x},\ldots,v_{Q,x}}_{k_{\mathrm{der}}\text{ arguments}}],
\qquad
0\le k_{\mathrm{der}}\le m,
\quad t_{\mathrm{seg}}\in\mathbb R,
\]
with order zero interpreted as the function value. Expanding the symmetric derivatives in coordinate directions and using the definition of \(\vartheta_Q\) yields
\[
\sum_{k_{\mathrm{der}}=0}^{m}
\frac{\varphi_{Q,x}^{(k_{\mathrm{der}})}(0)}{k_{\mathrm{der}}!}
=
\sum_{\alpha\in\mathcal A_{\mathrm{poly}}}
\frac{D^\alpha g_{\mathrm{ext}}(z_{Q,0})}{\alpha!}
\prod_{i=1}^{d}(v_{Q,x,i})^{\alpha_i}
=\Pi_Q(x),
\qquad
\varphi_{Q,x}(1)=\mu_{1,P}(x).
\]
Taylor's theorem with the order-\(m\) Lagrange remainder gives a point satisfying
\[
\exists\,t_{Q,x}\in(0,1):\qquad
\varphi_{Q,x}(1)
-
\sum_{k_{\mathrm{der}}=0}^{m-1}
\frac{\varphi_{Q,x}^{(k_{\mathrm{der}})}(0)}{k_{\mathrm{der}}!}
=
\frac{\varphi_{Q,x}^{(m)}(t_{Q,x})}{m!}.
\]
Choose such a point and define the degree-\(m\) remainder by
\[
\begin{aligned}
R_{\mathrm{Taylor},Q,x}
&:=\mu_{1,P}(x)-\Pi_Q(x)\\
&=\frac{\varphi_{Q,x}^{(m)}(t_{Q,x})
-\varphi_{Q,x}^{(m)}(0)}{m!}\\
&=\frac1{m!}
\left(
\mathrm D^m g_{\mathrm{ext}}(z_{Q,0}+t_{Q,x}v_{Q,x})
-\mathrm D^m g_{\mathrm{ext}}(z_{Q,0})
\right)
[\underbrace{v_{Q,x},\ldots,v_{Q,x}}_{m\text{ arguments}}].
\end{aligned}
\]
The displayed Hölder modulus for the extended function therefore implies
\[
\begin{aligned}
|R_{\mathrm{Taylor},Q,x}|
&\le
\frac{K_{\mathrm{ext}}K_{\mathrm{aff}}L}{m!}
\|t_{Q,x}v_{Q,x}\|_\infty^s
\|v_{Q,x}\|_\infty^m\\
&=
\frac{K_{\mathrm{ext}}K_{\mathrm{aff}}L}{m!}
 t_{Q,x}^{s}\|v_{Q,x}\|_\infty^{m+s}\\
&\le
\frac{K_{\mathrm{ext}}K_{\mathrm{aff}}L}{m!}
\|v_{Q,x}\|_\infty^\beta,
\end{aligned}
\]
since \(0<t_{Q,x}<1\), \(s>0\), and \(m+s=\beta\). This bound also applies when \(v_{Q,x}=0\). Define
\[
C_{\mathrm{Taylor}}:=\frac1{m!}\ge0.
\]
Because \(\|v_{Q,x}\|_\infty\le2h\), we obtain, uniformly in the expansion centre,
\[
\begin{aligned}
|\mu_{1,P}(x)-\Pi_Q(x)|
&=
\left|
g_{\mathrm{ext}}\!\left(
T_{\mathrm{aff}}(a_Q)+2h\frac{x-a_Q}{h}
\right)
-
U\!\left(\frac{x-a_Q}{h}\right)^\top\vartheta_Q
\right|
\\
&\le
C_{\mathrm{Taylor}}K_{\mathrm{ext}}K_{\mathrm{aff}}L(2h)^\beta
=
C_{\mathrm{bias}}Lh^\beta,
\end{aligned}
\]
where
\[
C_{\mathrm{bias}}
:=
C_{\mathrm{Taylor}}K_{\mathrm{ext}}K_{\mathrm{aff}}2^\beta.
\]
The normalized coordinates belong to \([0,1]^d\), and the affine image of \(x\) belongs to \([-1,1]^d\), which justify the equality and the approximation. This construction also covers integer \(\beta\), for which \(s=1\).
% lean: CausalSmith.Stat.WeakOverlap.exists_uniform_holderBall_local_monoVec_approx

\emph{2. Conditional bias and supremum loss.}
Apply \cref{obj:lem:selected-mesh}. Denote its uniform mesh and coefficient constants by \(K_{\mathrm{mesh}},K_{\mathrm{coef}}>0\), and its sample-size threshold by \(n_{\mathrm{env}}\). Enlarge these constants and the threshold as required by the constructions below, taking \(n_{\mathrm{env}}\ge2\). Define
\[
n_0:=\max\{n_{\mathrm{env}},1\},
\qquad
\mathcal G_n:=\sigma((X_i,A_i)_{i=1}^n).
\]
For \(n\ge n_0\), \(\gamma\in\Gamma\), and \(P\in\mathcal P_\gamma\), choose the simultaneous treated-count event underlying the mesh and coefficient bounds as follows. For all candidate microcells, define
\[
q_{Q\ell}:=P\{A=1,\ X\in E_{Q\ell}\},
\qquad h\in\mathcal H_n,\quad Q\in\mathcal Q_h,\quad 1\le\ell\le J,
\]
\[
M_{\mathrm{count},n}
:=\sum_{h\in\mathcal H_n}J|\mathcal Q_h|,
\qquad
s_{\mathrm{count},n}
:=\log\!\left(2(M_{\mathrm{count},n}+1)n^2\right)>0,
\]
and take
\[
\mathcal E_n
:=
\bigcap_{h\in\mathcal H_n}
\bigcap_{Q\in\mathcal Q_h}
\bigcap_{\ell=1}^J
\left\{
\left|N_{Q\ell}-nq_{Q\ell}\right|
<\frac{nq_{Q\ell}}2+4s_{\mathrm{count},n}
\right\}.
\]
The Bernstein bound for each treated count and the finite union bound give
\[
\begin{aligned}
P^{\otimes n}(\mathcal E_n^c)
&\le 2M_{\mathrm{count},n}e^{-s_{\mathrm{count},n}}
\\
&=\frac{M_{\mathrm{count},n}}{M_{\mathrm{count},n}+1}n^{-2}
\le n^{-2}.
\end{aligned}
\]
We establish the selected-mesh and treated-count bounds on this particular intersection. Let \(\eta\) be the reference microcube width in \cref{obj:def:template}. Each microcell at mesh \(h\) has volume \((\eta h)^d\), so \cref{obj:ass:density} and the measurable-set bound in \cref{obj:lem:ordered-mass} give
\[
\begin{aligned}
P(X\in E_{Q\ell})&\ge c_f(\eta h)^d,\\
q_{Q\ell}
&\ge \frac{\gamma-1}{\gamma}C^{-1/(\gamma-1)}
\bigl(c_f(\eta h)^d\bigr)^{\gamma/(\gamma-1)}
=\kappa_{\mathrm{cell}}(\gamma)h^{D_\gamma},
\end{aligned}
\]
where
\[
\kappa_{\mathrm{cell}}(\gamma)
:=\frac{\gamma-1}{\gamma}C^{-1/(\gamma-1)}
(c_f\eta^d)^{\gamma/(\gamma-1)},
\qquad
\kappa_{\mathrm{cell},\Gamma}
:=\min_{\gamma'\in\Gamma}\kappa_{\mathrm{cell}}(\gamma')>0.
\]
The positivity of the minimum follows from continuity and positivity of the displayed coefficient on the compact interval \(\Gamma\). Choose a fixed enlargement constant satisfying
\[
A_{\mathrm{or}}\ge1,
\qquad
4\le\kappa_{\mathrm{cell},\Gamma}
A_{\mathrm{or}}^{2\beta+D_{\gamma_{\max}}}.
\]
Such a choice exists because \(2\beta+D_{\gamma_{\max}}>0\). Since
\[
D_\gamma=d+\frac d{\gamma-1},
\qquad
D_{\gamma_{\max}}\le D_\gamma\le D_{\gamma_{\min}},
\qquad
h_{\gamma_{\max},n}\le h_{\gamma,n}\le h_{\gamma_{\min},n},
\]
we have
\[
4\le\kappa_{\mathrm{cell},\Gamma}
A_{\mathrm{or}}^{2\beta+D_\gamma}
\quad(\gamma\in\Gamma).
\]

The candidate collection in \cref{obj:synth_3} has largest index
\[
j_{\mathrm{cand},n}
:=\left\lfloor\frac{\log_2 n}{2\beta+d}\right\rfloor,
\qquad
\mathcal H_n=\{2^{-j}:0\le j\le j_{\mathrm{cand},n}\}.
\]
The floor inequality gives
\[
2^{-j_{\mathrm{cand},n}}
\le2n^{-1/(2\beta+d)}.
\]
Because \(d>0\) and \(\gamma_{\max}>1\),
\[
D_{\gamma_{\max}}>d,
\qquad
\frac{2n^{-1/(2\beta+d)}}{A_{\mathrm{or}}h_{\gamma_{\max},n}}
=\frac2{A_{\mathrm{or}}}
 n^{-1/(2\beta+d)+1/(2\beta+D_{\gamma_{\max}})}
\longrightarrow0,
\qquad
A_{\mathrm{or}}h_{\gamma_{\min},n}\longrightarrow0.
\]
After enlarging \(n_{\mathrm{env}}\), these relations ensure, uniformly over \(\gamma\in\Gamma\), that
\[
2^{-j_{\mathrm{cand},n}}
\le A_{\mathrm{or}}h_{\gamma,n}\le1.
\]
Define an oracle candidate by
\[
j_{\mathrm{or}}(n,\gamma)
:=\max\{j\in\{0,\ldots,j_{\mathrm{cand},n}\}:
2^{-j}\ge A_{\mathrm{or}}h_{\gamma,n}\},
\qquad
h_{\mathrm{or},n,\gamma}:=2^{-j_{\mathrm{or}}(n,\gamma)}.
\]
The set defining the maximum contains zero. If its maximum is smaller than \(j_{\mathrm{cand},n}\), the next dyadic width is smaller than \(A_{\mathrm{or}}h_{\gamma,n}\); if the maximum equals \(j_{\mathrm{cand},n}\), the preceding lower bound on the finest width forces equality with the target width. Thus
\[
h_{\mathrm{or},n,\gamma}\in\mathcal H_n,
\qquad
A_{\mathrm{or}}h_{\gamma,n}
\le h_{\mathrm{or},n,\gamma}
\le2A_{\mathrm{or}}h_{\gamma,n}.
\]

We next verify the logarithmic budget. The template and dyadic partition have
\[
J=(m+1)^d,
\qquad
|\mathcal Q_{2^{-j}}|=2^{jd},
\qquad
M_{\mathrm{count},n}
=J\sum_{j=0}^{j_{\mathrm{cand},n}}2^{jd}.
\]
Since \(2\beta+d\ge1\), the definition of \(j_{\mathrm{cand},n}\) gives
\[
j_{\mathrm{cand},n}\le2^{j_{\mathrm{cand},n}}\le n,
\qquad
M_{\mathrm{count},n}\le(n+1)n^dJ.
\]
For \(n\ge2\), using \(n+1\le2n\) and \(n^{d+1}\ge1\), we obtain
\[
2(M_{\mathrm{count},n}+1)n^2
\le8(J+1)n^{d+3},
\qquad
s_{\mathrm{count},n}
\le(d+3)\log n+\log(8(J+1)).
\]
Define
\[
U_{\mathrm{count}}:=24(d+3),
\qquad
V_{\mathrm{count}}:=24\log(8(J+1)),
\qquad
\rho_{\mathrm{budget}}
:=\frac{2\beta}{2\beta+D_{\gamma_{\min}}}>0.
\]
The elementary bound
\[
\log n\le\frac2{\rho_{\mathrm{budget}}}
 n^{\rho_{\mathrm{budget}}/2}
\]
and divergence of positive powers allow a further uniform enlargement of \(n_{\mathrm{env}}\) such that
\[
2U_{\mathrm{count}}\frac2{\rho_{\mathrm{budget}}}
\le(2A_{\mathrm{or}})^{-2\beta}
 n^{\rho_{\mathrm{budget}}/2},
\qquad
2V_{\mathrm{count}}
\le(2A_{\mathrm{or}})^{-2\beta}n^{\rho_{\mathrm{budget}}}.
\]
Consequently,
\[
\begin{aligned}
24s_{\mathrm{count},n}
&\le U_{\mathrm{count}}\log n+V_{\mathrm{count}}\\
&\le(2A_{\mathrm{or}})^{-2\beta}n^{\rho_{\mathrm{budget}}}
=(2A_{\mathrm{or}}h_{\gamma_{\min},n})^{-2\beta}\\
&\le(2A_{\mathrm{or}}h_{\gamma,n})^{-2\beta}
\le h_{\mathrm{or},n,\gamma}^{-2\beta}.
\end{aligned}
\]
The population-mass bound also gives, for every microcell of this oracle candidate,
\[
\begin{aligned}
nq_{Q\ell}
&\ge n\kappa_{\mathrm{cell},\Gamma}
 h_{\mathrm{or},n,\gamma}^{D_\gamma}\\
&=\kappa_{\mathrm{cell},\Gamma}
 n h_{\mathrm{or},n,\gamma}^{2\beta+D_\gamma}
 h_{\mathrm{or},n,\gamma}^{-2\beta}\\
&\ge\kappa_{\mathrm{cell},\Gamma}
 A_{\mathrm{or}}^{2\beta+D_\gamma}
 h_{\mathrm{or},n,\gamma}^{-2\beta}
\ge4h_{\mathrm{or},n,\gamma}^{-2\beta},
\end{aligned}
\]
where \(nh_{\gamma,n}^{2\beta+D_\gamma}=1\).

On \(\mathcal E_n\), the defining deviation inequality therefore yields, at the oracle candidate,
\[
N_{Q\ell}
>\frac{nq_{Q\ell}}2-4s_{\mathrm{count},n}
\ge2h_{\mathrm{or},n,\gamma}^{-2\beta}
 -4s_{\mathrm{count},n}
\ge h_{\mathrm{or},n,\gamma}^{-2\beta}.
\]
Thus this candidate is feasible. The smallest-feasible-mesh rule of \cref{obj:def:mesh-selector} gives
\[
0<\widehat h\le h_{\mathrm{or},n,\gamma}
\le2A_{\mathrm{or}}h_{\gamma,n},
\qquad
24s_{\mathrm{count},n}
\le h_{\mathrm{or},n,\gamma}^{-2\beta}
\le\widehat h^{-2\beta}.
\]
For a microcell of the selected mesh, feasibility and the same deviation inequality imply
\[
N_{Q\ell}\ge\widehat h^{-2\beta}
\ge24s_{\mathrm{count},n},
\qquad
\frac{nq_{Q\ell}}2
<N_{Q\ell}+4s_{\mathrm{count},n}
\le\frac76N_{Q\ell}.
\]
In particular,
\[
N_{Q\ell}\ge\frac{nq_{Q\ell}}5
\quad(Q\in\mathcal Q_{\widehat h},\ 1\le\ell\le J).
\]
Taking \(K_{\mathrm{mesh}}\ge2A_{\mathrm{or}}\), we have established on the displayed count intersection that
\[
0<\widehat h\le K_{\mathrm{mesh}}h_{\gamma,n},
\qquad
N_{Q\ell}\ge\frac{nq_{Q\ell}}5.
\]
% lean: CausalSmith.Stat.WeakOverlap.selectedCountEvent_uniform_design_saturated
% lean: CausalSmith.Stat.WeakOverlap.selectedCount_deterministic_bridge
Every count in the defining finite intersection is measurable with respect to \(\mathcal G_n\), and every population mass is fixed under \(P\). Hence \(\mathcal E_n\in\mathcal G_n\). Conditioning therefore fixes the sampled design, the selected mesh, its treated counts, and membership in \(\mathcal E_n\).

First, the response bound holds everywhere on the cube. Define
\[
g_{\mathrm{clip}}(x)
:=
\operatorname{clip}_{[-B,B]}(\mu_{1,P}(x)),
\qquad x\in\mathcal X.
\]
The conditional-mean bound in \cref{obj:ass:bounded-potential-outcomes} and identification give
\[
g_{\mathrm{clip}}(x)
=
\mu_{1,P}(x)
=
\mathbb E_P[Y\mid A=1,X=x]
\quad\text{for }P_X\text{-almost every }x.
\]
Both functions are continuous on the cube. The uniqueness conclusion of \cref{obj:lem:causal-identification} therefore gives
\[
g_{\mathrm{clip}}(x)=\mu_{1,P}(x),
\qquad
|\mu_{1,P}(x)|\le B
\quad(x\in\mathcal X).
\]

On \(\mathcal E_n\), set
\[
h:=\widehat h,
\qquad
\bar c_Q:=\mathbb E[\widehat c_Q\mid\mathcal G_n],
\qquad Q\in\mathcal Q_h.
\]
At each fixed design, this conditional expectation evaluates the coefficients at that design's selected mesh. Feasibility gives positive treated counts. For each treated observation used in the fit on \(Q\), define
\[
U_{i,Q}:=U\!\left(\frac{X_i-a_Q}{h}\right),
\qquad
R_{\mathrm{res},i,Q}
:=\mu_{1,P}(X_i)-U_{i,Q}^{\top}\vartheta_Q.
\]
The local polynomial approximation gives
\[
|R_{\mathrm{res},i,Q}|\le C_{\mathrm{bias}}Lh^\beta.
\]
Identification and the equal-cell normal equations yield, almost surely,
\[
\bar c_Q-\vartheta_Q
=
\widehat G_Q^{-1}
\frac1J\sum_{\ell=1}^J\frac1{N_{Q\ell}}
\sum_{\substack{i:A_i=1\\X_i\in E_{Q\ell}}}
U_{i,Q}R_{\mathrm{res},i,Q}.
\]
Every monomial coordinate has absolute value at most one on the normalized cube, so
\[
\|U_{i,Q}\|_2\le\sqrt{p_{\mathrm{coef}}}.
\]
Consequently, \cref{obj:lem:template-conditioning} implies
\[
\|\bar c_Q-\vartheta_Q\|_\infty
\le
\|\bar c_Q-\vartheta_Q\|_2
\le
\frac{2\sqrt{p_{\mathrm{coef}}}}{\lambda_0}
C_{\mathrm{bias}}Lh^\beta.
\]
Define the fixed constants
\[
K_{\mathrm{inv}}
:=\frac{2\sqrt{p_{\mathrm{coef}}}}{\lambda_0},
\qquad
K_{\mathrm{bias}}
:=p_{\mathrm{coef}}K_{\mathrm{inv}}C_{\mathrm{bias}}L
+C_{\mathrm{bias}}L.
\]

For the selected design, define the maximum coefficient fluctuation by
\[
S_{\mathrm{coef}}
:=
\max_{Q\in\mathcal Q_h}
\|\widehat c_Q-\bar c_Q\|_\infty.
\]
This maximum uses the selected mesh fixed by the conditioning design. Since
\[
|U(z)^\top(\widehat c_Q-\vartheta_Q)|
\le
p_{\mathrm{coef}}
\bigl(S_{\mathrm{coef}}+K_{\mathrm{inv}}C_{\mathrm{bias}}Lh^\beta\bigr),
\qquad z\in[0,1]^d,
\]
and clipping contracts distance to every value in \([-B,B]\), we obtain
\[
\sup_{x\in\mathcal X}
|\widehat\mu_n(x)-\mu_{1,P}(x)|
\le
p_{\mathrm{coef}}S_{\mathrm{coef}}+K_{\mathrm{bias}}h^\beta.
\]
We derive the conditional coefficient envelope on this same count event. For the selected design, define
\[
q_Q:=\min_{1\le\ell\le J}q_{Q\ell},
\qquad
\rho_\gamma:=\frac1{\gamma-1},
\]
and order the selected cubes so that
\[
\{Q_{(k)}:1\le k\le h^{-d}\}=\mathcal Q_h,
\qquad
q_{(k)}:=q_{Q_{(k)}},
\qquad
q_{(1)}\le\cdots\le q_{(h^{-d})}.
\]
Here \(h^{-d}=|\mathcal Q_h|\ge1\). The ordered-mass conclusion of \cref{obj:lem:ordered-mass} supplies a fixed constant satisfying
\[
\kappa_\Gamma>0,
\qquad
q_{(k)}\ge\kappa_\Gamma h^{D_\gamma}k^{\rho_\gamma}
\quad(1\le k\le h^{-d}).
\]

Define the treated residuals and coefficient weights by
\[
\epsilon_{\mathrm{tr},i}
:=\mathbf1_{\{A_i=1\}}\bigl(Y_i-\mu_{1,P}(X_i)\bigr),
\]
\[
w_{\mathrm{fit},i,Q,\alpha}
:=\frac1J\sum_{\ell=1}^J
\frac{\mathbf1_{\{A_i=1,\ X_i\in E_{Q\ell}\}}}{N_{Q\ell}}
\bigl(\widehat G_Q^{-1}U_{i,Q}\bigr)_\alpha,
\qquad \alpha\in\mathcal A_{\mathrm{poly}}.
\]
These weights are evaluated only at the positive feasible mesh on \(\mathcal E_n\), where every denominator is positive. Conditional on \(\mathcal G_n\), the residuals are independent by iid sampling, and identification together with \cref{obj:ass:bounded-potential-outcomes} gives, almost surely,
\[
\mathbb E[\epsilon_{\mathrm{tr},i}\mid\mathcal G_n]=0,
\qquad
\mathbb E[e^{t\epsilon_{\mathrm{tr},i}}\mid\mathcal G_n]
\le e^{B^2t^2/2}
\quad(t\in\mathbb R).
\]
For an untreated observation the residual is zero. The equal-cell coefficient formula consequently gives
\[
\widehat c_{Q,\alpha}-\bar c_{Q,\alpha}
=\sum_{i=1}^n w_{\mathrm{fit},i,Q,\alpha}
\epsilon_{\mathrm{tr},i}.
\]
The inverse bound in \cref{obj:lem:template-conditioning} and \(\|U_{i,Q}\|_2\le\sqrt{p_{\mathrm{coef}}}\) imply
\[
\left|\frac1J
(\widehat G_Q^{-1}U_{i,Q})_\alpha\right|
\le\frac{K_{\mathrm{inv}}}{J}.
\]
Cauchy--Schwarz over the \(J\) microcells, followed by summation over their treated observations, gives
\[
\begin{aligned}
\sum_{i=1}^n w_{\mathrm{fit},i,Q,\alpha}^2
&\le J\left(\frac{K_{\mathrm{inv}}}{J}\right)^2
\sum_{\ell=1}^J\frac1{N_{Q\ell}}\\
&\le K_{\mathrm{inv}}^2
\min\left(h^{2\beta},\frac5{nq_Q}\right).
\end{aligned}
\]
The last line uses both \(N_{Q\ell}\ge h^{-2\beta}\) and the established bound \(N_{Q\ell}\ge nq_{Q\ell}/5\). Multiplying the conditional exponential moments of the independent weighted residuals yields
\[
\mathbb E\!\left[
 e^{t(\widehat c_{Q,\alpha}-\bar c_{Q,\alpha})}
 \,\middle|\,\mathcal G_n\right]
\le\exp\!\left\{
\frac{B^2t^2}{2}\sum_{i=1}^n
w_{\mathrm{fit},i,Q,\alpha}^2\right\}.
\]
Define
\[
K_{\mathrm{mgf}}
:=B^2K_{\mathrm{inv}}^2\max\{1,5/\kappa_\Gamma\}>0.
\]
The ordered-mass bound and
\(\min\{u,cv\}\le\max\{1,c\}\min\{u,v\}\) for nonnegative \(u,v,c\) now give
\[
\mathbb E\!\left[
 e^{t(\widehat c_{Q_{(k)},\alpha}-\bar c_{Q_{(k)},\alpha})}
 \,\middle|\,\mathcal G_n\right]
\le\exp\!\left\{
\frac{K_{\mathrm{mgf}}t^2}{2}
\min\left(h^{2\beta},n^{-1}h^{-D_\gamma}k^{-\rho_\gamma}\right)
\right\}.
\]
All these exponential moments are finite; their two-sided finiteness also gives conditional integrability of the finite coefficient maxima.
% lean: CausalSmith.Stat.WeakOverlap.selectedMesh_saturated_rank_mgf

To apply the maximal inequality in \cref{obj:lem:selected-mesh} with a constant uniform in \(\gamma\), define
\[
\rho_{\min}:=\frac1{\gamma_{\max}-1},
\qquad
\rho_{\max}:=\frac1{\gamma_{\min}-1},
\qquad
R_{\mathrm{max}}:=\frac{\rho_{\max}}{\rho_{\min}},
\]
\[
a_{\mathrm{max}}:=\sqrt{K_{\mathrm{mgf}}}\,h^\beta,
\qquad
b_{\mathrm{max}}:=\sqrt{K_{\mathrm{mgf}}}\,n^{-1/2}h^{-D_\gamma/2}.
\]
Then
\[
0<\rho_{\min}\le\rho_\gamma\le\rho_{\max},
\qquad R_{\mathrm{max}}\ge1,
\qquad
 a_{\mathrm{max}}^2=K_{\mathrm{mgf}}h^{2\beta},
\qquad
 b_{\mathrm{max}}^2=K_{\mathrm{mgf}}n^{-1}h^{-D_\gamma}.
\]
Since \(k^{-\rho_\gamma}\le k^{-\rho_{\min}}\) for \(k\ge1\), the same moment bound satisfies the maximal inequality at the fixed exponent \(\rho_{\min}\). Let \(K_{\min,\mathrm{max}}>0\) be its constant, and define
\[
M_{\mathrm{max}}:=K_{\min,\mathrm{max}}\sqrt{R_{\mathrm{max}}}.
\]
The logarithmic factors satisfy
\[
\begin{aligned}
1+\max\left\{0,\log\left(
 (b_{\mathrm{max}}/a_{\mathrm{max}})^{2/\rho_{\min}}
\right)\right\}
\le R_{\mathrm{max}}\left[
1+\max\left\{0,\log\left(
 (b_{\mathrm{max}}/a_{\mathrm{max}})^{2/\rho_\gamma}
\right)\right\}\right].
\end{aligned}
\]
Indeed, when \(b_{\mathrm{max}}/a_{\mathrm{max}}\ge1\), expand the logarithms and use \(\rho_\gamma\le\rho_{\max}\); when this ratio is smaller than one, both positive-part logarithms vanish. Applying the scalar maximal inequality to each coefficient and using
\[
S_{\mathrm{coef}}
\le\sum_{\alpha\in\mathcal A_{\mathrm{poly}}}
\max_{Q\in\mathcal Q_h}
|\widehat c_{Q,\alpha}-\bar c_{Q,\alpha}|,
\]
we obtain
\[
\mathbb E[S_{\mathrm{coef}}\mid\mathcal G_n]
\le p_{\mathrm{coef}}M_{\mathrm{max}}a_{\mathrm{max}}
\sqrt{1+\max\left\{0,\log\left(
(b_{\mathrm{max}}/a_{\mathrm{max}})^{2/\rho_\gamma}
\right)\right\}}.
\]

The oracle-mesh identity gives
\[
n^{-1/2}=h_{\gamma,n}^{\beta+D_\gamma/2},
\qquad
\frac{b_{\mathrm{max}}}{a_{\mathrm{max}}}
=\left(\frac{h_{\gamma,n}}h\right)^{\beta+D_\gamma/2}.
\]
Define
\[
\zeta_{\mathrm{log},\gamma}:=2\beta(\gamma-1)+d\gamma,
\qquad
\zeta_{\mathrm{log},\Gamma}
:=\max\{1,2\beta(\gamma_{\max}-1)+d\gamma_{\max}\}.
\]
Then
\[
\log\left((b_{\mathrm{max}}/a_{\mathrm{max}})^{2/\rho_\gamma}\right)
=\zeta_{\mathrm{log},\gamma}\log(h_{\gamma,n}/h),
\qquad
0<\zeta_{\mathrm{log},\gamma}\le\zeta_{\mathrm{log},\Gamma},
\]
and hence
\[
\mathbb E[S_{\mathrm{coef}}\mid\mathcal G_n]
\le p_{\mathrm{coef}}M_{\mathrm{max}}
\sqrt{K_{\mathrm{mgf}}\zeta_{\mathrm{log},\Gamma}}\,
 h^\beta\sqrt{1+\max\{0,\log(h_{\gamma,n}/h)\}}.
\]
To absorb the remaining logarithm, define
\[
z_{\mathrm{mesh}}:=h/h_{\gamma,n},
\qquad
0<z_{\mathrm{mesh}}\le K_{\mathrm{mesh}}.
\]
If \(z_{\mathrm{mesh}}\le1\), then
\[
\log(1/z_{\mathrm{mesh}})\le1/z_{\mathrm{mesh}}-1,
\qquad
\sqrt{1+\log(1/z_{\mathrm{mesh}})}\le1/z_{\mathrm{mesh}},
\qquad
z_{\mathrm{mesh}}^\beta\le z_{\mathrm{mesh}},
\]
so their product is at most one. If \(z_{\mathrm{mesh}}>1\), the positive-part logarithm vanishes and \(z_{\mathrm{mesh}}^\beta\le K_{\mathrm{mesh}}^\beta\). Thus
\[
h^\beta\sqrt{1+\max\{0,\log(h_{\gamma,n}/h)\}}
\le\max\{1,K_{\mathrm{mesh}}^\beta\}r_{\gamma,n}.
\]
Choose the initially enlarged coefficient constant to satisfy
\[
K_{\mathrm{coef}}\ge
p_{\mathrm{coef}}M_{\mathrm{max}}
\sqrt{K_{\mathrm{mgf}}\zeta_{\mathrm{log},\Gamma}}
\max\{1,K_{\mathrm{mesh}}^\beta\}.
\]
We have therefore established, almost surely on the explicitly chosen count event,
\[
\mathbb E[S_{\mathrm{coef}}\mid\mathcal G_n]
\le K_{\mathrm{coef}}r_{\gamma,n}
\quad\text{on }\mathcal E_n.
\]
% lean: CausalSmith.Stat.WeakOverlap.selectedMesh_saturated_conditional_envelope
Also,
\[
h^\beta
\le K_{\mathrm{mesh}}^\beta r_{\gamma,n}
\le \max\{1,K_{\mathrm{mesh}}^\beta\}r_{\gamma,n}.
\]
Thus, defining
\[
K_{\mathrm{cond}}
:=
p_{\mathrm{coef}}K_{\mathrm{coef}}
+
K_{\mathrm{bias}}\max\{1,K_{\mathrm{mesh}}^\beta\}
+1,
\]
we have \(K_{\mathrm{cond}}>0\) and
\[
\mathbb E\!\left[
\sup_{x\in\mathcal X}
|\widehat\mu_n(x)-\mu_{1,P}(x)|
\,\middle|\,\mathcal G_n
\right]
\le K_{\mathrm{cond}}r_{\gamma,n}
\quad\text{almost surely on }\mathcal E_n.
\]
The finite mesh selection and polynomial fits are measurable, and the conditional integrability needed here is supplied by \cref{obj:lem:selected-mesh}.
% lean: CausalSmith.Stat.WeakOverlap.selectedMesh_saturated_conditional_cubeLoss

\emph{3. Uniform expected risk.}
For every sample, including samples with \(\widehat h=0\), the estimator in \cref{obj:def:estimator} takes values in \([-B,B]\). Define, on the entire sample space,
\[
Z_{\mathrm{loss}}
:=
\sup_{x\in\mathcal X}
|\widehat\mu_n(x)-\mu_{1,P}(x)|,
\qquad
F_{\mathrm{loss}}
:=
\mathbb E[Z_{\mathrm{loss}}\mid\mathcal G_n].
\]
The preceding response bound gives
\[
0\le Z_{\mathrm{loss}}\le2B,
\qquad
0\le F_{\mathrm{loss}}\le2B.
\]
The supremum loss is measurable and bounded. The tower property and the bounds on the event and its complement therefore give
\[
\begin{aligned}
\mathbb E_{P^{\otimes n}}Z_{\mathrm{loss}}
&=\mathbb E_{P^{\otimes n}}F_{\mathrm{loss}}
\\
&\le
K_{\mathrm{cond}}r_{\gamma,n}
P^{\otimes n}(\mathcal E_n)
+
2B P^{\otimes n}(\mathcal E_n^c)
\\
&\le K_{\mathrm{cond}}r_{\gamma,n}+2Bn^{-2}.
\end{aligned}
\]
For every \(n\ge1\), \(\beta>0\), and \(\gamma>1\),
\[
D_\gamma\ge0,
\qquad
2\beta+D_\gamma>0,
\qquad
\frac{\beta}{2\beta+D_\gamma}\le1,
\]
so monotonicity of powers with base \(n\ge1\) gives
\[
n^{-2}
\le n^{-\beta/(2\beta+D_\gamma)}
=r_{\gamma,n}.
\]
Set
\[
K:=K_{\mathrm{cond}}+2B.
\]
Then \(K>0\), and
\[
\mathbb E_{P^{\otimes n}}
\left[
\sup_{x\in\mathcal X}
|\widehat\mu_n(x)-\mu_{1,P}(x)|
\right]
\le Kr_{\gamma,n}
\]
uniformly for \(n\ge n_0\), \(\gamma\in\Gamma\), and \(P\in\mathcal P_\gamma\). The count rule and clipped fit in \cref{obj:def:mesh-selector,obj:def:estimator} use the stated inputs \(n,d,\beta,B,\mathcal D_n\).
% lean: CausalSmith.Stat.WeakOverlap.estimatorSupRisk_uniform_upper

\emph{4. The bounded-subclass lower bound.}
Fix \(\gamma>1\) and \(x_0\in\mathcal X\). Define the subclass and its pointwise minimax risk by
\[
\mathcal P_\gamma^{\{0,B\}}
:=
\left\{
P\in\mathcal P_\gamma:
P\text{ admits a causal completion with }
Y(0),Y(1)\in\{0,B\}\text{ almost surely}
\right\},
\]
\[
R_{\mathrm{bd}}(n,\gamma,x_0)
:=
\inf_T\sup_{P\in\mathcal P_\gamma^{\{0,B\}}}
\mathbb E_{P^{\otimes n}}
|T(\mathcal D_n,x_0)-\mu_{1,P}(x_0)|.
\]
The infimum is over sample-measurable scalar estimators.

Apply \cref{obj:lem:bounded-lower-pair}. It supplies fixed \(a,c,c_0>0\) such that, for every \(n\ge1\), the pair from \cref{obj:def:bounded-lower-pair}, with
\[
h:=ch_{\gamma,n},
\]
belongs to \(\mathcal P_\gamma^{\{0,B\}}\) and satisfies
\[
P_{0,h}\ll P_{1,h},
\qquad
\log\frac{dP_{0,h}}{dP_{1,h}}\in L^1(P_{0,h}),
\qquad
n\,\operatorname{KL}(P_{0,h},P_{1,h})\le\frac18.
\]
Define
\[
\mathbb P_{j,n}:=P_{j,h}^{\otimes n},
\qquad
\tau_{j,n}:=\mu_{1,P_{j,h}}(x_0)
\quad(j\in\{0,1\}),
\qquad
\delta_n:=c_0r_{\gamma,n}.
\]
The separation supplied by \cref{obj:lem:bounded-lower-pair} is
\[
0<\delta_n\le|\tau_{1,n}-\tau_{0,n}|.
\]

For an arbitrary sample-measurable estimator \(T\), define its clipped version by
\[
T_{\mathrm{clip}}(\mathcal D_n)
:=
\operatorname{clip}_{[-B,B]}
\bigl(T(\mathcal D_n,x_0)\bigr).
\]
The response bound established above gives, for \(j\in\{0,1\}\),
\[
|\tau_{j,n}|\le B,
\qquad
|T_{\mathrm{clip}}-\tau_{j,n}|
\le |T(\mathcal D_n,x_0)-\tau_{j,n}|,
\qquad
|T_{\mathrm{clip}}-\tau_{j,n}|\le2B.
\]
Thus the clipped losses are integrable under both sample laws.

Product entropy tensorization and Pinsker's inequality give
\[
\operatorname{KL}(\mathbb P_{0,n},\mathbb P_{1,n})
\le n\,\operatorname{KL}(P_{0,h},P_{1,h})
\le\frac18\le\frac12,
\]
\[
\|\mathbb P_{0,n}-\mathbb P_{1,n}\|_{\mathrm{TV}}
\le
\sqrt{\frac{\operatorname{KL}(\mathbb P_{0,n},\mathbb P_{1,n})}{2}}
\le\frac12.
\]
Define the two large-error events by
\[
\mathcal A_{j,n}
:=
\left\{
\mathcal D_n:
|T_{\mathrm{clip}}(\mathcal D_n)-\tau_{j,n}|
\ge\frac{\delta_n}{2}
\right\},
\qquad j\in\{0,1\}.
\]
The triangle inequality and target separation imply
\[
\mathcal A_{0,n}^c\subseteq\mathcal A_{1,n}.
\]
Hence
\[
\begin{aligned}
\mathbb P_{0,n}(\mathcal A_{0,n})
+\mathbb P_{1,n}(\mathcal A_{1,n})
&\ge
1-\|\mathbb P_{0,n}-\mathbb P_{1,n}\|_{\mathrm{TV}},
\\
\max_{j\in\{0,1\}}\mathbb P_{j,n}(\mathcal A_{j,n})
&\ge
\frac{1-
\sqrt{\operatorname{KL}(\mathbb P_{0,n},\mathbb P_{1,n})/2}}{2}
\ge\frac14.
\end{aligned}
\]
For each \(j\), integration over \(\mathcal A_{j,n}\) yields
\[
\mathbb E_{\mathbb P_{j,n}}
|T_{\mathrm{clip}}-\tau_{j,n}|
\ge\frac{\delta_n}{2}\mathbb P_{j,n}(\mathcal A_{j,n}).
\]
Combining the two inequalities and the clipping contraction gives
\[
\begin{aligned}
\max_{j\in\{0,1\}}
\mathbb E_{\mathbb P_{j,n}}
|T(\mathcal D_n,x_0)-\tau_{j,n}|
&\ge
\max_{j\in\{0,1\}}
\mathbb E_{\mathbb P_{j,n}}
|T_{\mathrm{clip}}-\tau_{j,n}|
\\
&\ge\frac{\delta_n}{8}
=\frac{c_0}{8}r_{\gamma,n}.
\end{aligned}
\]
Both laws belong to the restricted subclass. Taking the infimum over \(T\), and defining
\[
c_\gamma:=\frac{c_0}{8}>0,
\]
proves
\[
c_\gamma r_{\gamma,n}
\le R_{\mathrm{bd}}(n,\gamma,x_0)
\qquad(n\ge1).
\]
% lean: CausalSmith.Stat.WeakOverlap.boundedPointwiseRisk_oracle_lower

\emph{5. Constants for all positive sample sizes.}
For the fixed \(\gamma\), define a local adaptation interval by
\[
\gamma_{\mathrm{lo}}:=\frac{\gamma+1}{2},
\qquad
\gamma_{\mathrm{hi}}:=\gamma+1.
\]
Since \(\gamma>1\),
\[
1<\gamma_{\mathrm{lo}}<\gamma_{\mathrm{hi}},
\qquad
\gamma\in[\gamma_{\mathrm{lo}},\gamma_{\mathrm{hi}}].
\]
Applying the uniform estimator bound to this interval supplies \(K_{\mathrm{loc}}>0\) and \(n_{\mathrm{loc}}\in\mathbb N\) such that
\[
\sup_{P\in\mathcal P_\gamma}
\mathbb E_{P^{\otimes n}}
\left[
\sup_{x\in\mathcal X}
|\widehat\mu_n(x)-\mu_{1,P}(x)|
\right]
\le K_{\mathrm{loc}}r_{\gamma,n}
\quad(n\ge\max\{n_{\mathrm{loc}},1\}).
\]
Define
\[
N_\gamma:=\max\{n_{\mathrm{loc}},1\},
\qquad
q_{\mathrm{small}}:=N_\gamma^{-2},
\qquad
K_{\mathrm{small}}:=\frac{2B}{q_{\mathrm{small}}},
\]
\[
K_\gamma
:=
\max\{c_\gamma,\max\{K_{\mathrm{loc}},K_{\mathrm{small}}\}\}.
\]
These constants are finite, \(q_{\mathrm{small}}>0\), \(K_{\mathrm{small}}>0\), and
\[
0<c_\gamma\le K_\gamma,
\qquad
K_{\mathrm{loc}}\le K_\gamma,
\qquad
K_{\mathrm{small}}\le K_\gamma.
\]

Fix \(n\ge1\). If \(n\ge n_{\mathrm{loc}}\), the measurable estimator \(\widehat\mu_n\) is an admissible witness in the supremum minimax infimum, so
\[
R_\infty(n,\gamma)
\le
\sup_{P\in\mathcal P_\gamma}
\mathbb E_{P^{\otimes n}}
\left[
\sup_{x\in\mathcal X}
|\widehat\mu_n(x)-\mu_{1,P}(x)|
\right]
\le K_{\mathrm{loc}}r_{\gamma,n}
\le K_\gamma r_{\gamma,n}.
\]
If \(n<n_{\mathrm{loc}}\), then \(1\le n\le N_\gamma\). Monotonicity of the inverse square and \(n^{-2}\le r_{\gamma,n}\) give
\[
q_{\mathrm{small}}
=N_\gamma^{-2}
\le n^{-2}
\le r_{\gamma,n}.
\]
Therefore
\[
2B
=K_{\mathrm{small}}q_{\mathrm{small}}
\le K_{\mathrm{small}}r_{\gamma,n}
\le K_\gamma r_{\gamma,n}.
\]
The same clipped estimator has loss at most \(2B\) for every sample, so
\[
R_\infty(n,\gamma)
\le
\sup_{P\in\mathcal P_\gamma}
\mathbb E_{P^{\otimes n}}
\left[
\sup_{x\in\mathcal X}
|\widehat\mu_n(x)-\mu_{1,P}(x)|
\right]
\le2B
\le K_\gamma r_{\gamma,n}.
\]
Thus the upper bound holds for every \(n\ge1\).
% lean: CausalSmith.Stat.WeakOverlap.adaptive_minimax_rate

\emph{6. Comparison of the risks.}
The inclusion
\[
\mathcal P_\gamma^{\{0,B\}}\subseteq\mathcal P_\gamma
\]
implies, for each scalar estimator \(T\),
\[
\sup_{P\in\mathcal P_\gamma^{\{0,B\}}}
\mathbb E_{P^{\otimes n}}
|T(\mathcal D_n,x_0)-\mu_{1,P}(x_0)|
\le
\sup_{P\in\mathcal P_\gamma}
\mathbb E_{P^{\otimes n}}
|T(\mathcal D_n,x_0)-\mu_{1,P}(x_0)|.
\]
Taking infima over the same estimators gives
\[
R_{\mathrm{bd}}(n,\gamma,x_0)
\le R_{\mathrm{pt}}(n,\gamma,x_0).
\]
For any admissible function-valued estimator \(T\), define its evaluation by
\[
T_{x_0}(\mathcal D_n):=T(\mathcal D_n)(x_0).
\]
This is a measurable scalar estimator, and \(x_0\in\mathcal X\) gives
\[
|T_{x_0}(\mathcal D_n)-\mu_{1,P}(x_0)|
\le
\sup_{x\in\mathcal X}
|T(\mathcal D_n)(x)-\mu_{1,P}(x)|.
\]
Taking expectations, worst-case risks, and then the infimum over function-valued estimators yields
\[
R_{\mathrm{pt}}(n,\gamma,x_0)\le R_\infty(n,\gamma).
\]
Combining the bounded-subclass lower bound, the upper bound for every positive sample size, and these comparisons proves
\[
c_\gamma r_{\gamma,n}
\le R_{\mathrm{bd}}(n,\gamma,x_0)
\le R_{\mathrm{pt}}(n,\gamma,x_0)
\le R_\infty(n,\gamma)
\le K_\gamma r_{\gamma,n},
\qquad n\ge1.
\]
% lean: CausalSmith.Stat.WeakOverlap.boundedPointwiseRisk_le_pointwiseRisk
% lean: CausalSmith.Stat.WeakOverlap.pointwiseRisk_le_supremumRisk
\end{proof}

\begin{proof}[Proof of \cref{obj:prop:strict-enlargement}]
Fix an admissible strip exponent \(a\), and use the construction and selected propensity in \cref{obj:def:thin-strip} throughout.

\begin{enumerate}
\item Let the completed-data law and its treated response be
\[
P_{\mathrm{comp}}
:=
\mathcal L(X,A,Y(0),Y(1),Y),
\qquad
\mu_1(x):=Bp_0,\quad x\in\mathcal X.
\]
The construction gives
\[
0<p_0=\min\left\{\frac14,\frac{L}{4B}\right\}\le\frac14,
\qquad
Bp_0\le B\frac{L}{4B}=\frac L4.
\]
The radial term and strip increment are nonnegative, and truncation at one therefore gives
\[
0\le e_\star(x)\le1,\qquad x\in\mathcal X.
\]
Thus all conditional Bernoulli distributions are probability distributions. Integrating their product against the uniform covariate distribution shows that \(P_{\mathrm{comp}}\) and its observed-data marginal \(P_{\mathrm{strip}}\) are probability laws, with
\[
(P_{\mathrm{comp}})_X=(P_{\mathrm{strip}})_X
=\operatorname{Leb}_d\!\restriction_{\mathcal X}.
\]

The conditional distribution of \(Y(1)\) is the constant scaled Bernoulli distribution, so
\[
\mathbb E_{P_{\mathrm{comp}}}[Y(1)\mid X]
=(1-p_0)\,0+p_0B
=Bp_0
=\mu_1(X)
\quad\text{almost surely}.
\]
For the derivative order \(m\) in \cref{obj:def:holder}, the constant function satisfies
\[
\max_{|\alpha|\le m}\|D^\alpha\mu_1\|_\infty
=Bp_0\le\frac L4\le L,
\qquad
\max_{|\alpha|=m}\sup_{x\ne x'}
\frac{|D^\alpha\mu_1(x)-D^\alpha\mu_1(x')|}
{\|x-x'\|_\infty^{\beta-m}}
=0.
\]
Its derivatives have continuous boundary traces. Hence
\[
\mu_1\in\mathcal H^\beta(L),
\]
as defined in \cref{obj:def:holder}, and it supplies the response required by \cref{obj:ass:holder-response}.
% lean: CausalSmith.Stat.WeakOverlap.completionBind_holderResponse_constant_equal_arms

\item We verify the remaining restrictions in \cref{obj:def:model}. The uniform marginal gives
\[
P_{\mathrm{strip}}(X\in\mathcal X)=1,
\qquad
f(x)=1\ge c_f
\quad\text{for Lebesgue-almost every }x\in\mathcal X,
\]
which supplies \cref{obj:ass:density}. The completed construction gives
\[
Y=AY(1)+(1-A)Y(0)
\quad\text{almost surely},
\]
and, for almost every \(x\in\mathcal X\),
\[
\mathcal L\bigl((Y(0),Y(1)),A\mid X=x\bigr)
=
\mathcal L\bigl(Y(0),Y(1)\mid X=x\bigr)
\otimes\operatorname{Bernoulli}(e_\star(x)).
\]
These identities verify \cref{obj:ass:consistency,obj:ass:exchangeability}, and the treatment marginal gives
\[
P_{\mathrm{comp}}(A=1\mid X=x)=e_\star(x)
\quad\text{for almost every }x.
\]

For the overlap calculation, define the radial propensity and radial exponent by
\[
q_{\mathrm{rad}}(x)
:=F_R(R(x))^{1/(\gamma-1)},\quad x\in\mathcal X,
\qquad
\delta_{\mathrm{strip}}:=\frac d{\gamma-1}>0.
\]
The centered cube of radius \(r\) has side length \(2r\), whence
\[
F_R(r)=(2r)^d,\qquad 0\le r\le\frac12.
\]
In particular \(0\le q_{\mathrm{rad}}(x)\le1\), so the nonnegative strip increment gives
\[
q_{\mathrm{rad}}(x)\le e_\star(x),
\qquad
\{x\in\mathcal X:e_\star(x)\le t\}
\subseteq
\{x\in\mathcal X:F_R(R(x))\le t^{\gamma-1}\}.
\]

To bound the last event over its whole range, for \(0\le\tau_{\mathrm{cdf}}<1\) define
\[
r_{\mathrm{cdf}}(\tau_{\mathrm{cdf}})
:=\frac{\tau_{\mathrm{cdf}}^{1/d}}2.
\]
Then \(0\le r_{\mathrm{cdf}}(\tau_{\mathrm{cdf}})\le1/2\), and monotonicity of the \(d\)-th power on nonnegative reals gives
\[
\{x\in\mathcal X:F_R(R(x))\le\tau_{\mathrm{cdf}}\}
\subseteq
\{x\in\mathcal X:R(x)\le r_{\mathrm{cdf}}(\tau_{\mathrm{cdf}})\}.
\]
Consequently,
\[
P_{\mathrm{strip}}\{F_R(R(X))\le\tau_{\mathrm{cdf}}\}
\le
\bigl(2r_{\mathrm{cdf}}(\tau_{\mathrm{cdf}})\bigr)^d
=\tau_{\mathrm{cdf}}.
\]
This includes \(\tau_{\mathrm{cdf}}=0\); at \(\tau_{\mathrm{cdf}}=1\), the same upper bound follows from total probability one. Taking \(\tau_{\mathrm{cdf}}=t^{\gamma-1}\) therefore yields
\[
P_{\mathrm{strip}}\{e_\star(X)\le t\}
\le t^{\gamma-1}
\le Ct^{\gamma-1},
\qquad 0\le t\le1.
\]
This verifies \cref{obj:ass:global-tail}. At \(t=0\), the sharp bound and nonnegativity of \(e_\star\) also give
\[
P_{\mathrm{strip}}\{e_\star(X)=0\}=0,
\qquad
e_\star(X)>0\quad\text{almost surely}.
\]

The equal-arm product construction identifies the treated conditional distribution wherever treatment has positive conditional probability:
\[
\mathcal L_{P_{\mathrm{strip}}}(Y\mid A=1,X=x)
=\mathcal L\bigl(B\operatorname{Bernoulli}(p_0)\bigr),
\qquad
\mathbb E_{P_{\mathrm{strip}}}[Y\mid A=1,X=x]=Bp_0
\]
for \(P_X\)-almost every \(x\). Its support is contained in \([0,B]\). The bounded-range exponential-moment inequality, followed by enlargement of its exponent, gives
\[
\mathbb E_{P_{\mathrm{strip}}}
\!\left[\exp\{\lambda(Y-Bp_0)\}\mid A=1,X=x\right]
\le
\exp\left\{\frac{B^2\lambda^2}{8}\right\}
\le
\exp\left\{\frac{B^2\lambda^2}{2}\right\},
\qquad \lambda\in\mathbb R.
\]
Also \(|Bp_0|\le B\). For every \(\lambda\ne0\), it follows that
\[
\frac{
\sqrt{2\log\mathbb E_{P_{\mathrm{strip}}}
[\exp\{\lambda(Y-Bp_0)\}\mid A=1,X=x]}
}{|\lambda|}
\le B.
\]
The centered exponential moment is at least one by Jensen's inequality, so the square root is well defined. These estimates supply \cref{obj:ass:bounded-potential-outcomes}. Together with the smoothness already established and the stated parameter conditions, they prove
\[
P_{\mathrm{strip}}\in\mathcal P_\gamma.
\]
The uniform covariate marginal also supplies the asserted sharp bound
\[
P_{\mathrm{strip}}\{e_\star(X)\le t\}\le t^{\gamma-1},
\qquad t\in[0,1].
\]
% lean: CausalSmith.Stat.WeakOverlap.thinStrip_globalTailModel_of_holder
% lean: CausalSmith.Stat.WeakOverlap.stripPropensity_tail_of_uniform_covariateLaw

\item Fix arbitrary \(\rho,\nu,h_0>0\). For \(h>0\), define the centered Euclidean ball and local propensity supremum by
\[
\mathcal B_h:=\{x\in\mathbb R^d:\|x-x^\circ\|_2\le h\},
\qquad
M_{\mathrm{loc}}(h)
:=\sup_{x\in\mathcal B_h\cap\mathcal X}e_\star(x).
\]
For \(0<h\le1/2\), every coordinate displacement in \(\mathcal B_h\) has magnitude at most \(h\). On \(S\cap\mathcal B_h\), the strip inequality further gives
\[
|x_2-1/2|\le |x_1-1/2|^a\le h^a.
\]
Thus this intersection is contained in a rectangular box with second-coordinate side length \(2h^a\) and all other side lengths \(2h\). Uniformity implies
\[
P_{\mathrm{strip}}(X\in S\cap\mathcal B_h)
\le(2h^a)(2h)^{d-1}.
\]
For the denominator, the centered box with every coordinate displacement at most \(h/d\) lies in \(\mathcal X\) and in \(\mathcal B_h\), since
\[
\sum_{i=1}^d(x_i-1/2)^2
\le d(h/d)^2
=\frac{h^2}{d}\le h^2.
\]
Its volume therefore supplies
\[
P_{\mathrm{strip}}(X\in\mathcal B_h)
\ge(2h/d)^d>0.
\]
Dividing these bounds yields
\[
\frac{P_{\mathrm{strip}}(X\in S\cap\mathcal B_h)}
{P_{\mathrm{strip}}(X\in\mathcal B_h)}
\le
\frac{(2h^a)(2h)^{d-1}}{(2h/d)^d}
=d^d h^{a-1}.
\]
Because \(a-1>0\) and \(\delta_{\mathrm{strip}}>0\),
\[
d^d h^{a-1}\longrightarrow0,
\qquad
(2h)^{\delta_{\mathrm{strip}}}\longrightarrow0
\quad\text{as }h\downarrow0.
\]
Choose \(h\) satisfying
\[
0<h<\min\{h_0,1/2\},
\qquad
d^d h^{a-1}<\nu,
\qquad
(2h)^{\delta_{\mathrm{strip}}}<\rho/4.
\]

The center belongs to \(S\), and its radial term is zero. Hence
\[
e_\star(x^\circ)=\frac14,
\qquad
M_{\mathrm{loc}}(h)\ge\frac14.
\]
For \(x\in\mathcal B_h\setminus S\), the propensity equals its radial component and \(R(x)\le h\), giving
\[
e_\star(x)=q_{\mathrm{rad}}(x)
\le(2h)^{\delta_{\mathrm{strip}}}
<\rho/4
\le\rho M_{\mathrm{loc}}(h).
\]
Therefore
\[
\{x\in\mathcal B_h:e_\star(x)\ge\rho M_{\mathrm{loc}}(h)\}
\subseteq S\cap\mathcal B_h,
\]
and
\[
\frac{
P_{\mathrm{strip}}\{e_\star(X)\ge\rho M_{\mathrm{loc}}(h),
\ X\in\mathcal B_h\}
}{
P_{\mathrm{strip}}(X\in\mathcal B_h)
}
\le d^d h^{a-1}<\nu.
\]
This contradicts the numerical clause for the proposed triple. Since the triple was arbitrary, the numerical anti-concentration clause fails. The argument uses the pointwise bound \(e_\star\le1\), so it applies with propensity values equal to one admitted.
% lean: CausalSmith.Stat.WeakOverlap.thinStripLaw_not_dornA3AntiConcentration

\item The full singleton-family condition in \cref{obj:def:dorn-a3} entails its numerical clause with the same constants \(\rho,\nu,h_0\). Thus the preceding contradiction gives
\[
\neg\mathsf A_3^{\mathrm{Dorn}}(P_{\mathrm{strip}})
\]
for the selected propensity \(e_\star\).
% lean: CausalSmith.Stat.WeakOverlap.dornA3AntiConcentrationOnCube_of_dornA3OnCube

\item For the design limit, define, for every \(h>0\),
\[
\mathcal Q_{\mathrm{ctr},h}
:=\{x\in\mathbb R^d:R(x)\le h\},
\]
\[
\begin{aligned}
D_{\mathrm{strip}}(h)
&:=P_{\mathrm{strip}}\{A=1,\ X\in\mathcal Q_{\mathrm{ctr},h}\},\\
N_{\mathrm{strip}}(h)
&:=\int_{\{A=1,\ X\in\mathcal Q_{\mathrm{ctr},h}\}}
\left(\frac{X_2-1/2}{h}\right)^2\,dP_{\mathrm{strip}},\\
T_{\mathrm{strip}}(h)
&:=\frac{N_{\mathrm{strip}}(h)}{D_{\mathrm{strip}}(h)}.
\end{aligned}
\]
For this definition, a quotient with zero denominator is assigned value zero.

Introduce the positive constants and envelope
\[
b_{\mathrm{strip}}
:=\frac{1}{2(4d)^{a+d-1}},
\qquad
c_{\mathrm{rad}}:=2^{d+\delta_{\mathrm{strip}}},
\qquad
c_{\mathrm{strip}}:=2^d,
\]
\[
G_{\mathrm{strip}}(h)
:=
\frac{
c_{\mathrm{rad}}h^{d+\delta_{\mathrm{strip}}}
+c_{\mathrm{strip}}h^{d+3a-3}
}{
b_{\mathrm{strip}}h^{d+a-1}
},
\qquad h>0.
\]
Factoring powers of \(h\) gives
\[
G_{\mathrm{strip}}(h)
=
\frac{c_{\mathrm{rad}}}{b_{\mathrm{strip}}}
h^{1+\delta_{\mathrm{strip}}-a}
+
\frac{c_{\mathrm{strip}}}{b_{\mathrm{strip}}}
h^{2a-2}.
\]
The strip-exponent restrictions give
\[
1+\delta_{\mathrm{strip}}-a>0,
\qquad
2a-2>0,
\]
so both powers tend to zero and
\[
G_{\mathrm{strip}}(h)\longrightarrow0
\quad\text{as }h\downarrow0.
\]
% lean: CausalSmith.Stat.WeakOverlap.thinStrip_raw_second_envelope_tendsto_zero

\item We bound the moment by this envelope. Suppose \(0<h\le1/2\). On \(\mathcal Q_{\mathrm{ctr},h}\), the normalized second-coordinate square is at most one. If \(x\in S\cap\mathcal Q_{\mathrm{ctr},h}\), the strip inequality and \(e_\star(x)\le1\) give
\[
\left(\frac{x_2-1/2}{h}\right)^2e_\star(x)
\le\left(\frac{h^a}{h}\right)^2.
\]
If \(x\in\mathcal Q_{\mathrm{ctr},h}\setminus S\), the radial bound gives
\[
\left(\frac{x_2-1/2}{h}\right)^2e_\star(x)
\le e_\star(x)\le(2h)^{\delta_{\mathrm{strip}}}.
\]
Combining these two cases and integrating,
\[
\begin{aligned}
&\int_{\mathcal Q_{\mathrm{ctr},h}}
\left(\frac{x_2-1/2}{h}\right)^2e_\star(x)\,dx\\
&\qquad\le
(2h)^{\delta_{\mathrm{strip}}}
P_{\mathrm{strip}}(X\in\mathcal Q_{\mathrm{ctr},h})
+
\left(\frac{h^a}{h}\right)^2
P_{\mathrm{strip}}(X\in S\cap\mathcal Q_{\mathrm{ctr},h}).
\end{aligned}
\]
Here \(\mathcal Q_{\mathrm{ctr},h}\subseteq\mathcal X\), and its volume is \((2h)^d\). The same rectangular enclosure used above bounds the strip intersection, so
\[
\begin{aligned}
&\int_{\mathcal Q_{\mathrm{ctr},h}}
\left(\frac{x_2-1/2}{h}\right)^2e_\star(x)\,dx\\
&\qquad\le
(2h)^{\delta_{\mathrm{strip}}}(2h)^d
+\left(\frac{h^a}{h}\right)^2(2h^a)(2h)^{d-1}\\
&\qquad=
c_{\mathrm{rad}}h^{d+\delta_{\mathrm{strip}}}
+c_{\mathrm{strip}}h^{d+3a-3}.
\end{aligned}
\]

For the treated-mass lower bound, define
\[
r_{\mathrm{core}}(h):=\frac{h}{4d},
\]
\[
\mathcal K_h
:=
\left\{
x\in\mathbb R^d:
|x_2-1/2|\le r_{\mathrm{core}}(h)^a,\
r_{\mathrm{core}}(h)\le x_i-1/2\le2r_{\mathrm{core}}(h)
\ \text{for every }i\ne2
\right\}.
\]
Since \(0<r_{\mathrm{core}}(h)\le1/8\) and \(a>1\),
\[
r_{\mathrm{core}}(h)^a\le r_{\mathrm{core}}(h).
\]
The coordinate bounds place \(\mathcal K_h\) in \(\mathcal X\), and they imply
\[
|x_2-1/2|
\le r_{\mathrm{core}}(h)^a
\le |x_1-1/2|^a,
\]
\[
\|x-x^\circ\|_2^2
\le d\bigl(2r_{\mathrm{core}}(h)\bigr)^2
=\frac{h^2}{4d}\le h^2.
\]
Thus
\[
\mathcal K_h\subseteq S\cap\mathcal B_h
\subseteq S\cap\mathcal Q_{\mathrm{ctr},h}.
\]
Its second-coordinate side length is \(2r_{\mathrm{core}}(h)^a\), and each other side length is \(r_{\mathrm{core}}(h)\). Since \(e_\star\ge1/4\) on \(S\),
\[
\begin{aligned}
\int_{\mathcal Q_{\mathrm{ctr},h}}e_\star(x)\,dx
&\ge\int_{\mathcal B_h\cap\mathcal X}e_\star(x)\,dx\\
&\ge\frac14
\bigl(2r_{\mathrm{core}}(h)^a\bigr)
r_{\mathrm{core}}(h)^{d-1}\\
&=b_{\mathrm{strip}}h^{d+a-1}>0.
\end{aligned}
\]

Finally, conditioning on \(X\) and using
\(P_{\mathrm{strip}}(A=1\mid X)=e_\star(X)\) gives the numerator and denominator identities
\[
N_{\mathrm{strip}}(h)
=
\int_{\mathcal Q_{\mathrm{ctr},h}}
\left(\frac{x_2-1/2}{h}\right)^2e_\star(x)\,dx,
\qquad
D_{\mathrm{strip}}(h)
=
\int_{\mathcal Q_{\mathrm{ctr},h}}e_\star(x)\,dx.
\]
All integrands are nonnegative. Dividing the numerator upper bound by the positive denominator lower bound therefore yields
\[
0\le T_{\mathrm{strip}}(h)\le G_{\mathrm{strip}}(h),
\qquad 0<h\le1/2.
\]
The squeeze theorem proves
\[
\lim_{h\downarrow0}T_{\mathrm{strip}}(h)=0,
\]
which is exactly the asserted normalized treated-design limit.
% lean: CausalSmith.Stat.WeakOverlap.stripSecondMoment_tendsto_zero

\item The law \(P_{\mathrm{strip}}\) itself supplies the final existence claim: it belongs to \(\mathcal P_\gamma\), and its full singleton-family Dorn predicate fails for \(e_\star\).
% lean: CausalSmith.Stat.WeakOverlap.thinStrip_mem_model_not_dornA3
\end{enumerate}
\end{proof}

\begin{proof}[Proof of \cref{obj:prop:dorn-a3-free-corollary}]
\begin{enumerate}
\item \emph{The source mean is bounded throughout the cube.}
Fix \(\gamma\in\Gamma\) and a source tuple in \(\mathcal D_\gamma\). Write
\[
X^{\mathrm D}:=X,
\qquad
P_X^{\mathrm D}:=\mathcal L_P(X^{\mathrm D})
=2^{-d}\operatorname{Leb}\big|_{[-1,1]^d},
\qquad
\mu_{1,P}:=\mu_P,
\]
and let \(\mu_{0,P}\) be its selected control mean. Continuity of \(\mu_P\) on the compact source cube gives
\[
\int |\mu_P(z)|\,P_X^{\mathrm D}(dz)<\infty.
\]
For \(P_X^{\mathrm D}\)-almost every \(z\), the treated conditional distribution is
\(N(\mu_P(z),\sigma_{\min}^2)\). Its mean identity, the absolute-integral inequality, and Jensen's inequality give
\[
|\mu_P(z)|
\le
\int |y|\,N(\mu_P(z),\sigma_{\min}^2)(dy),
\]
and
\[
|\mu_P(z)|^q
\le
\left(\int |y|\,N(\mu_P(z),\sigma_{\min}^2)(dy)\right)^q
\le
\int |y|^q\,N(\mu_P(z),\sigma_{\min}^2)(dy)
\le M^q.
\]
Since \(q>0\) and \(M>0\), this implies \(|\mu_P|\le M\) almost everywhere.

To extend this inequality to the boundary as well, define
\[
v_P(z):=\max\{|\mu_P(z)|-M,0\},
\qquad z\in\mathbb R^d.
\]
The function \(v_P\) is continuous on the source cube and vanishes Lebesgue-almost everywhere there. Moreover,
\[
[-1,1]^d
=
\overline{\operatorname{int}([-1,1]^d)}.
\]
Continuity first forces \(v_P\) to vanish on the interior, since a positive value would give a neighborhood of positive Lebesgue measure on which it remains positive. Taking limits from the interior then gives
\[
v_P(z)=0,
\qquad
|\mu_P(z)|\le M
\qquad(z\in[-1,1]^d).
\]
% lean: CausalSmith.Stat.WeakOverlap.dorn_A3_free_corollary

\item \emph{Completion and comparison of the common radii.}
By \cref{obj:lem:fixed-cube-holder-completion}, there is a constant \(K_{d,\beta}>0\) such that, for every nonnegative amplitude bound \(M'\), nonnegative seminorm bound \(L_0'\), and function satisfying the corresponding source regularity and bounds,
\[
u\circ T
\in
\mathcal H^\beta\!\left(K_{d,\beta}(M'+L_0')\right).
\]
For a function admitting a nonnegative Hölder radius, write its set of admissible radii as
\[
\mathscr L_\beta(g)
:=
\{L_{\mathrm{adm}}\in[0,\infty):
g\in\mathcal H^\beta(L_{\mathrm{adm}})\},
\qquad
\|g\|_{\mathcal H^\beta}
=
\inf\mathscr L_\beta(g).
\]
Consequently,
\[
\|u\circ T\|_{\mathcal H^\beta}
\le K_{d,\beta}(M'+L_0').
\]
In particular, the two suprema in the proposition are bounded above:
\[
C_{d,\beta}\le K_{d,\beta},
\qquad
L_*\le K_{d,\beta}(M+L_0).
\]

Every admissible radius bounds the function itself. Taking the infimum of those bounds gives
\[
|g(x)|\le \|g\|_{\mathcal H^\beta},
\qquad x\in[0,1]^d.
\]
Constant functions satisfy the source seminorm restriction because their derivative differences vanish. Applying the preceding inequality to the constant functions \(1\) and \(M\), which admit radii by the completion result, yields
\[
1\le \|1\|_{\mathcal H^\beta},
\qquad
M\le \|M\|_{\mathcal H^\beta}.
\]
The constant function \(1\), with amplitude and seminorm bounds both equal to \(1\), is admissible in the definition of \(C_{d,\beta}\). The constant function \(M\) is admissible in the definition of \(L_*\). Hence
\[
0<\frac12\le C_{d,\beta}\le K_{d,\beta}<\infty,
\qquad
0<M\le L_*.
\]
Finally, every function contributing to \(L_*\) also contributes, after normalization, to \(C_{d,\beta}\):
\[
\frac{\|u\circ T\|_{\mathcal H^\beta}}{M+L_0}
\le C_{d,\beta}.
\]
Taking the supremum over these functions proves
\[
0<L_*\le C_{d,\beta}(M+L_0)<\infty.
\]
% lean: CausalSmith.Stat.WeakOverlap.dorn_A3_free_corollary

\item \emph{Gaussian source witnesses for the local-condition assertion.}
Suppose \(d\ge2\), and fix \(\gamma\in\Gamma\). Then \(\gamma>1\). Choose
\[
\delta_{\mathrm{strip}}:=\frac{d}{\gamma-1}>0,
\qquad
a:=1+\frac{\delta_{\mathrm{strip}}}{2},
\]
so that
\[
1<a<1+\frac{d}{\gamma-1}.
\]
For the fixed moment exponent, set
\[
I_q:=\int |y|^q\,N(0,1)(dy),
\qquad
M':=\max\{1,|I_q|\},
\qquad
C':=\left((3/4)^{\gamma-1}\right)^{-1},
\qquad
\sigma_{\min}':=1,
\qquad
L_0':=1.
\]
Gaussian moments are finite, so \(M'\) is positive and finite. Since \(M'\ge1\) and \(q>1\),
\[
I_q\le |I_q|\le M'\le (M')^q.
\]
All four chosen source constants are positive and finite.

Use the center, radial function, strip, and propensity from
\cref{obj:def:thin-strip}. For the ambient arguments needed below, the same coordinate formulas are
\[
x^\circ=(1/2,\ldots,1/2),
\qquad
R(x)=\|x-x^\circ\|_\infty
\quad(x\in\mathbb R^d),
\]
\[
F_R(t)=
\begin{cases}
0,&t<0,\\
\min\{1,(2t)^d\},&t\ge0,
\end{cases}
\qquad
\widetilde S:=
\{x\in\mathbb R^d:
|x_2-1/2|\le |x_1-1/2|^a\}.
\]
Here \(\widetilde S\cap[0,1]^d=S\). Define the capped propensity on all of \(\mathbb R^d\) by
\[
e_{\mathrm{cap}}(x)
:=
\min\left\{\frac34,\,
\min\left\{1,\,
F_R(R(x))^{1/(\gamma-1)}
+\frac14\mathbf1_{\widetilde S}(x)\right\}\right\}.
\]
Thus, on the unit cube,
\[
e_{\mathrm{cap}}(x)=\min\{3/4,e_\star(x)\}.
\]

Construct random variables with conditional product distribution
\[
X_{\mathrm{strip}}\sim\operatorname{Unif}([0,1]^d),
\]
\[
\mathcal L\!\left(
A_{\mathrm{strip}},Y_{\mathrm{strip}}(0),Y_{\mathrm{strip}}(1)
\mid X_{\mathrm{strip}}=x
\right)
=
\operatorname{Bernoulli}(e_{\mathrm{cap}}(x))
\otimes N(0,1)\otimes N(0,1).
\]
Set
\[
Y_{\mathrm{strip}}
:=
A_{\mathrm{strip}}Y_{\mathrm{strip}}(1)
+(1-A_{\mathrm{strip}})Y_{\mathrm{strip}}(0),
\qquad
X_{\mathrm G}^{\mathrm D}:=T(X_{\mathrm{strip}}),
\]
\[
P_{\mathrm G}:=
\mathcal L(X_{\mathrm G}^{\mathrm D},A_{\mathrm{strip}},Y_{\mathrm{strip}}),
\qquad
\mu_{1,\mathrm G}(z)=\mu_{0,\mathrm G}(z):=0,
\qquad
e_{\mathrm G}(z):=e_{\mathrm{cap}}(T^{-1}z)
\quad(z\in\mathbb R^d).
\]
These measurable product kernels define a probability law. The affine change of variables and the conditional product distribution give
\[
X_{\mathrm G}^{\mathrm D}\sim\operatorname{Unif}([-1,1]^d),
\qquad
e_{\mathrm G}(X_{\mathrm G}^{\mathrm D})
=
P_{\mathrm G}(A_{\mathrm{strip}}=1\mid X_{\mathrm G}^{\mathrm D}),
\]
and
\[
Y_{\mathrm{strip}}
\mid X_{\mathrm G}^{\mathrm D},A_{\mathrm{strip}}=b
\sim N(0,1),
\qquad b\in\{0,1\}.
\]
Since the radial term is positive away from \(x^\circ\), a set of full covariate probability,
\[
0<e_{\mathrm G}(X_{\mathrm G}^{\mathrm D})\le\frac34<1
\quad\text{almost surely}.
\]
The zero means satisfy the source smoothness restriction, and
\[
\operatorname{Var}_{P_{\mathrm G}}(\mu_{1,\mathrm G}(X_{\mathrm G}^{\mathrm D}))=0\le M',
\qquad
\operatorname{Var}_{P_{\mathrm G}}(Y_{\mathrm{strip}}\mid
X_{\mathrm G}^{\mathrm D},A_{\mathrm{strip}})=1=(\sigma_{\min}')^2.
\]
Together with the bound on \(I_q\), these establish the moment, variance, Gaussian, and smoothness restrictions.

The uniform-covariate propensity bound in
\cref{obj:prop:strict-enlargement} supplies
\[
\Pr\{e_\star(X_{\mathrm{strip}})\le t\}\le t^{\gamma-1},
\qquad 0\le t\le1.
\]
For \(0\le t<3/4\), capping leaves this sublevel event unchanged. Since \(C'\ge1\),
\[
\Pr\{e_{\mathrm{cap}}(X_{\mathrm{strip}})\le t\}
\le t^{\gamma-1}\le C't^{\gamma-1}.
\]
For \(3/4\le t\le1\), monotonicity of the positive power gives
\[
\Pr\{e_{\mathrm{cap}}(X_{\mathrm{strip}})\le t\}
\le1
=
C'(3/4)^{\gamma-1}
\le C't^{\gamma-1}.
\]
The affine transport therefore yields
\[
P_{\mathrm G}\{e_{\mathrm G}(X_{\mathrm G}^{\mathrm D})\le t\}
\le C't^{\gamma-1},
\qquad 0\le t\le1.
\]
Thus the tuple satisfies every retained source restriction with the chosen constants.

It remains to establish the failure of the numerical clause in
\cref{obj:def:dorn-a3} for this capped representative. Define the closed centered Euclidean ball by
\[
\mathcal B_h:=
\{x\in\mathbb R^d:\|x-x^\circ\|_2\le h\},
\qquad h>0.
\]
For \(0<h\le1/2\), the ball lies in the unit cube. On
\(S\cap\mathcal B_h\), the second coordinate has width at most \(2h^a\), and each other coordinate has width at most \(2h\). Conversely,
\[
\prod_{i=1}^d[1/2-h/d,1/2+h/d]\subseteq\mathcal B_h,
\]
because the squared Euclidean distance of any point in this box is at most
\(d(h/d)^2\le h^2\). These box bounds give
\[
\Pr\{X_{\mathrm{strip}}\in S\cap\mathcal B_h\}
\le (2h^a)(2h)^{d-1},
\]
\[
\Pr\{X_{\mathrm{strip}}\in\mathcal B_h\}
\ge (2h/d)^d>0,
\]
and hence
\[
\frac{\Pr\{X_{\mathrm{strip}}\in S\cap\mathcal B_h\}}
{\Pr\{X_{\mathrm{strip}}\in\mathcal B_h\}}
\le
\frac{(2h^a)(2h)^{d-1}}{(2h/d)^d}
=
d^d h^{a-1}.
\]

Consider any proposed constants \(\rho,\nu,h_0>0\), and set
\[
\rho_{\mathrm{cap}}:=\min\{\rho,2\}>0.
\]
Both \(h^{a-1}\) and \(((2h)^d)^{1/(\gamma-1)}\) tend to zero as \(h\downarrow0\). Choose \(h\) with
\[
0<h\le\min\{h_0,1/2\},
\qquad
d^d h^{a-1}\le\nu,
\qquad
((2h)^d)^{1/(\gamma-1)}
<
\frac{\rho_{\mathrm{cap}}}{4}.
\]
For every \(x\in\mathcal B_h\), \(R(x)\le h\), so
\[
F_R(R(x))^{1/(\gamma-1)}
\le ((2h)^d)^{1/(\gamma-1)}
<
\frac{\rho_{\mathrm{cap}}}{4}
\le\frac12.
\]
Consequently,
\[
e_\star(x)
\le F_R(R(x))^{1/(\gamma-1)}+\frac14
<\frac34,
\qquad
e_{\mathrm{cap}}(x)=e_\star(x)
\quad(x\in\mathcal B_h).
\]
At the center the propensity equals \(1/4\), giving the pointwise bound
\[
e_{\mathrm{cap}}(x^\circ)=\frac14,
\qquad
\sup_{x\in\mathcal B_h\cap[0,1]^d}e_{\mathrm{cap}}(x)\ge\frac14.
\]
For \(x\in\mathcal B_h\setminus S\), the strip term vanishes, and
\[
e_{\mathrm{cap}}(x)
=
F_R(R(x))^{1/(\gamma-1)}
<\frac{\rho}{4}.
\]
It follows that
\[
\left\{
x\in\mathcal B_h:
e_{\mathrm{cap}}(x)\ge
\rho\sup_{x'\in\mathcal B_h\cap[0,1]^d}e_{\mathrm{cap}}(x')
\right\}
\subseteq S\cap\mathcal B_h.
\]
The resulting local mass ratio is therefore at most
\(d^d h^{a-1}\le\nu\), contradicting the required strict inequality.

Finally, the affine identities
\[
\|T(x)-T(x^\circ)\|_2=2\|x-x^\circ\|_2,
\qquad
T(x^\circ)=\mathbf0,
\qquad
e_{\mathrm G}(T(x))=e_{\mathrm{cap}}(x)
\]
preserve the pointwise suprema and the numerator and denominator probabilities in these ratios. A source-cube local condition with radius bound \(h_0\) would thus give a unit-cube condition with radius bound \(h_0/2\), which the preceding argument contradicts. Hence
\[
\neg\,\mathsf A_3^{\mathrm{Dorn}}
\bigl(\{P_{\mathrm G}\},(e_{\mathrm G})\bigr).
\]
The construction selects its constants separately for the chosen \(\gamma\).
% lean: CausalSmith.Stat.WeakOverlap.gaussianThinStripOriginal_not_dornA3

\item \emph{Transport and causal completion of each source tuple.}
Return to an arbitrary tuple in \(\mathcal D_\gamma\), with \(\gamma\in\Gamma\). Define globally measurable versions by
\[
\mu_{b,P}^{\mathrm{meas}}(z):=
\begin{cases}
\mu_{b,P}(z),&z\in[-1,1]^d,\\
0,&z\notin[-1,1]^d,
\end{cases}
\qquad
b\in\{0,1\},\ z\in\mathbb R^d,
\]
\[
e_P^{\mathrm{cl}}(z):=
\max\{0,\min\{1,e_P(z)\}\},
\qquad z\in\mathbb R^d.
\]
Continuity on the closed source cube makes the mean extensions measurable. The clamped propensity is measurable and takes values in \([0,1]\). Since the source covariate is supported on the cube and the selected propensity lies in \((0,1)\) almost surely,
\[
\mu_{b,P}^{\mathrm{meas}}=\mu_{b,P},
\qquad
e_P^{\mathrm{cl}}=e_P
\quad P_X^{\mathrm D}\text{-almost everywhere}.
\]

Construct a completion law \(\Pi_P\) by retaining covariate marginal
\(P_X^{\mathrm D}\) and using the conditional product distribution
\[
\mathcal L_{\Pi_P}(A,Y(0),Y(1)\mid X^{\mathrm D}=z)
=
\operatorname{Bernoulli}(e_P^{\mathrm{cl}}(z))
\otimes
N(\mu_{0,P}^{\mathrm{meas}}(z),\sigma_{\min}^2)
\otimes
N(\mu_{1,P}^{\mathrm{meas}}(z),\sigma_{\min}^2),
\]
with
\[
Y:=AY(1)+(1-A)Y(0).
\]
All factors are probability kernels, so this defines a probability law.

To identify its observed marginal, take any bounded measurable function
\[
\varphi:\mathbb R^d\times\{0,1\}\times\mathbb R\longrightarrow\mathbb R.
\]
Conditioning first on the covariate and treatment and then on the covariate alone gives
\[
\begin{aligned}
\int \varphi\,dP
={}&
\int\Bigg[
(1-e_P^{\mathrm{cl}}(z))
\int\varphi(z,0,y)\,
N(\mu_{0,P}^{\mathrm{meas}}(z),\sigma_{\min}^2)(dy)
\\
&\hspace{34mm}
+
e_P^{\mathrm{cl}}(z)
\int\varphi(z,1,y)\,
N(\mu_{1,P}^{\mathrm{meas}}(z),\sigma_{\min}^2)(dy)
\Bigg]P_X^{\mathrm D}(dz)
\\
={}&
\mathbb E_{\Pi_P}\varphi(X^{\mathrm D},A,Y).
\end{aligned}
\]
Here the selected propensity and both Gaussian arm kernels identify the two conditional distributions. Thus
\[
\mathcal L_{\Pi_P}(X^{\mathrm D},A,Y)=P.
\]

Define the transported covariate, completion, and observed law by
\[
X^T:=T^{-1}(X^{\mathrm D}),
\]
\[
\Pi_P^T:=
\mathcal L_{\Pi_P}(X^T,A,Y(0),Y(1),Y),
\qquad
Q_P:=
\mathcal L_P(T^{-1}(X^{\mathrm D}),A,Y).
\]
The observed marginal of \(\Pi_P^T\) is exactly \(Q_P\). The source tail inequality at \(t=1\) also gives
\[
1=P\{e_P(X^{\mathrm D})\le1\}\le C.
\]
Together with \(B_*>0\), \(L_*>0\), and \(\gamma>1\), this verifies the required parameter domain.

The affine change of variables has Jacobian \(2^d\), and hence
\[
X^T\sim\operatorname{Unif}([0,1]^d),
\qquad
f_{Q_P}(x)=1
\quad\text{Lebesgue-almost everywhere on }[0,1]^d.
\]
The transported selected propensity is
\[
e_P^T(x):=e_P(T(x)),
\qquad x\in\mathbb R^d,
\]
and satisfies
\[
Q_P\{e_P^T(X^T)\le t\}
=
P\{e_P(X^{\mathrm D})\le t\}
\le Ct^{\gamma-1},
\qquad 0\le t\le1.
\]
The treated conditional kernel under \(Q_P\) is
\(N(\mu_P(T(x)),\sigma_{\min}^2)\) almost everywhere. The displayed bound \(|\mu_P(z)|\le M\) on the source cube gives
\[
|\mu_P(T(x))|\le M\le B_*,
\qquad x\in[0,1]^d.
\]
Translation of a Gaussian by its mean produces a centered Gaussian. Its moment-generating function therefore gives, for every \(\lambda\in\mathbb R\),
\[
\mathbb E_{Q_P}\!\left[
\exp\{\lambda(Y-\mu_P(T(X^T)))\}
\mid X^T,A=1
\right]
=
\exp\!\left(\frac{\lambda^2\sigma_{\min}^2}{2}\right)
\le
\exp\!\left(\frac{\lambda^2B_*^2}{2}\right).
\]
These inequalities verify \cref{obj:ass:bounded-potential-outcomes}.

The treated potential outcome is integrable, by integrability of the source mean and the Gaussian residual. Its Gaussian conditional distribution yields
\[
\mathbb E_{\Pi_P}[Y(1)\mid X^{\mathrm D}]
=
\mu_{1,P}^{\mathrm{meas}}(X^{\mathrm D}).
\]
Since \(T\) is a measurable bijection with measurable inverse, conditioning on \(X^T\) gives
\[
\mathbb E_{\Pi_P^T}[Y(1)\mid X^T]
=
\mu_{1,P}^{\mathrm{meas}}(T(X^T))
=
\mu_P(T(X^T))
\quad\text{almost surely}.
\]

We next verify smoothness at the precise radius \(L_*\). Completion supplies at least one admissible radius for \(\mu_P\circ T\), and its defining supremum gives
\[
\|\mu_P\circ T\|_{\mathcal H^\beta}\le L_*.
\]
For completeness, the infimum defining the full radius is attained whenever its admissible-radius set is nonempty. Indeed, regularity follows from any one admissible radius, and every derivative bound satisfies
\[
|D^\alpha g(x)|
\le
\inf\mathscr L_\beta(g)
=
\|g\|_{\mathcal H^\beta},
\qquad |\alpha|\le m.
\]
For distinct \(x,x'\), every admissible radius bounds the quotient, so
\[
\frac{|D^\alpha g(x)-D^\alpha g(x')|}
{\|x-x'\|_\infty^{\beta-m}}
\le
\inf\mathscr L_\beta(g)
=
\|g\|_{\mathcal H^\beta},
\qquad |\alpha|=m.
\]
For \(x=x'\), both sides of the corresponding modulus inequality have zero left-hand side and zero distance factor; it follows directly from any admissible radius. Thus
\[
g\in\mathcal H^\beta(\|g\|_{\mathcal H^\beta}).
\]
Increasing a radius preserves its defining inequalities. Applying this to the transported response proves
\[
\mu_P\circ T\in\mathcal H^\beta(L_*).
\]
The measurable mean extension agrees with the selected source mean throughout the source cube, so its transported version agrees with \(\mu_P\circ T\) throughout the unit cube. Their intrinsic derivative traces and modulus restrictions therefore coincide. Since \(X^T\) is supported there, this also preserves the conditional-mean identity.

Finally, the product construction and its affine transport give
\[
Y=AY(1)+(1-A)Y(0)
\quad\text{almost surely},
\qquad
(Y(0),Y(1))\perp A\mid X^T.
\]
These are \cref{obj:ass:consistency,obj:ass:exchangeability}. Together with the density, tail, response-bound, and smoothness calculations, they verify every restriction in \cref{obj:def:model}. Consequently,
\[
Q_P\in\mathcal P_\gamma
\quad\text{with}\quad
(B,L,C,c_f)=(B_*,L_*,C,1),
\]
with the exhibited causal completion \(\Pi_P^T\).
% lean: CausalSmith.Stat.WeakOverlap.dorn_A3_free_corollary

\item \emph{A common expected-supremum bound.}
Define the enlarged tail constant by
\[
C_{\mathrm{up}}:=\max\{C,1\}.
\]
Increasing the tail constant preserves the tail restriction, since
\[
Ct^{\gamma-1}\le C_{\mathrm{up}}t^{\gamma-1},
\qquad 0\le t\le1.
\]
Thus every transported tuple belongs to the model with parameters
\((d,\beta,B_*,L_*,C_{\mathrm{up}},1)\). These parameters satisfy the hypotheses of \cref{obj:thm:adaptive-minimax}.

To state the sampling identity explicitly, define, for \(n\in\mathbb N\),
\[
\Omega_n:=
(\mathbb R^d\times\{0,1\}\times\mathbb R)^n,
\]
\[
\mathcal T_n\bigl((z_i,a_i,y_i)_{i=1}^n\bigr)
:=
(T^{-1}z_i,a_i,y_i)_{i=1}^n,
\qquad
\mathcal T_n:\Omega_n\to\Omega_n.
\]
For \(n=0\), this is the identity on the singleton empty sample. For every source tuple, define the transported-sample loss by
\[
Z_{n,P}(\omega)
:=
\sup_{x\in[0,1]^d}
|\widehat\mu_n(\omega,x)-\mu_P(T(x))|,
\qquad \omega\in\Omega_n,
\]
where \(\widehat\mu_n\) is the estimator of \cref{obj:def:estimator} with parameters \(\beta,B_*\).

By \cref{obj:lem:causal-identification}, the continuous treated response associated with \(Q_P\) agrees throughout the cube with the response of the displayed completion:
\[
\mu_{1,Q_P}(x)=\mu_P(T(x)),
\qquad x\in[0,1]^d.
\]
Hence the estimator risk supplied by \cref{obj:thm:adaptive-minimax} is exactly
\[
\mathbb E_{Q_P^{\otimes n}}Z_{n,P}
=
\mathbb E_{P^{\otimes n}}
\bigl[Z_{n,P}(\mathcal T_n(\mathcal D_n))\bigr].
\]
That result supplies \(K>0\) and a common integer \(n_0\ge1\) such that
\[
\mathbb E_{Q_P^{\otimes n}}Z_{n,P}
\le K r_{\gamma,n},
\qquad
n\ge n_0,\quad \gamma\in\Gamma,\quad P\in\mathcal D_\gamma.
\]
The constants are common because the model parameters and adaptation interval are fixed.
% lean: CausalSmith.Stat.WeakOverlap.adaptive_minimax_rate

\item \emph{Restriction to source families.}
The source setting permits an arbitrary nonempty family whose members satisfy the retained restrictions with common constants
\citep[Section 2, Assumptions 1--2, pp.~3--7]{Dorn2026GlobalOverlap}.
Therefore every such family satisfies
\[
\mathfrak P\ne\varnothing,
\qquad
\mathfrak P\subseteq\mathcal D_\gamma.
\]
Taking the supremum of the preceding memberwise estimate over this family gives, with the same \(K,n_0\),
\[
\sup_{P\in\mathfrak P}
\mathbb E_{P^{\otimes n}}
\bigl[Z_{n,P}(\mathcal T_n(\mathcal D_n))\bigr]
\le K r_{\gamma,n},
\qquad n\ge n_0.
\]
This is the asserted family expected-supremum bound.
% lean: CausalSmith.Stat.WeakOverlap.dorn_A3_free_corollary

\item \emph{The probability bound on the source cube.}
Fix \(\varepsilon>0\), and choose
\[
c(\varepsilon):=\frac K\varepsilon>0,
\qquad
n_{0,\varepsilon}:=\max\{n_0,1\}.
\]
For every \(n\ge n_{0,\varepsilon}\), the oracle scale is positive:
\[
r_{\gamma,n}
=
n^{-\beta/(2\beta+d\gamma/(\gamma-1))}>0.
\]
For each partition cell \(Q\) at each fixed dyadic mesh, choose a countable subset dense in the relative topology:
\[
D_Q\subseteq Q,
\qquad
\overline{D_Q}^{\,Q}=Q.
\]
On a sample with positive selected mesh, the fitted function restricted to each selected cell is a clipped polynomial. Together with continuity of the transported response, this gives
\[
\sup_{x\in Q}
|\widehat\mu_n(\omega,x)-\mu_P(T(x))|
=
\sup_{x\in D_Q}
|\widehat\mu_n(\omega,x)-\mu_P(T(x))|,
\qquad Q\in\mathcal Q_{\widehat h(\omega)}.
\]
For each fixed mesh and cell, the fitted value at a fixed point is sample-measurable, so the right-hand side is a countable supremum of measurable functions. Taking the maximum over the finitely many cells preserves measurability. The mesh selector is measurable and takes values in a countable set of dyadic meshes together with zero, so selecting among these measurable losses also preserves measurability. On the zero-mesh branch,
\[
Z_{n,P}(\omega)
=
\sup_{x\in[0,1]^d}|\mu_P(T(x))|,
\qquad \widehat h(\omega)=0,
\]
which is deterministic. Thus \(Z_{n,P}\) is measurable.

For every transported sample \(\omega\) and every source-cube point \(z\),
\[
\left|
\widehat\mu_n(\omega,T^{-1}z)-\mu_P(z)
\right|
=
\left|
\widehat\mu_n(\omega,T^{-1}z)
-\mu_P(T(T^{-1}z))
\right|
\le Z_{n,P}(\omega).
\]
The source covariate law is supported on that cube, so
\[
\|\widehat\mu_n(\omega,\cdot)\circ T^{-1}-\mu_P\|_{L^\infty(P)}
\le Z_{n,P}(\omega).
\]
Markov's inequality and the common expectation bound now give
\[
\begin{aligned}
&Q_P^{\otimes n}\!\left\{
\|\widehat\mu_n\circ T^{-1}-\mu_P\|_{L^\infty(P)}
>c(\varepsilon)r_{\gamma,n}
\right\}
\\
&\qquad\le
Q_P^{\otimes n}\{Z_{n,P}>c(\varepsilon)r_{\gamma,n}\}
\le
\frac{\mathbb E_{Q_P^{\otimes n}}Z_{n,P}}
{(K/\varepsilon)r_{\gamma,n}}
\le\varepsilon.
\end{aligned}
\]
The sampling identity for \(\mathcal T_n\) expresses the same probability under \(P^{\otimes n}\) when the estimator is computed from the transported sample. Every bound is simultaneous over \(\gamma\in\Gamma\) and source tuples in \(\mathcal D_\gamma\).
% lean: CausalSmith.Stat.WeakOverlap.measurableLoss_probability_le

\item \emph{Families with a common regression curve.}
Let \(\mathfrak P\) be the stated nonempty family with common treated curve \(\mu^\circ\). For any \(n\in\mathbb N\) and evaluation point \(x_0\), define the constant sample estimator
\[
\widehat\theta_{\mathrm{known},n}(\omega):=\mu^\circ(x_0),
\qquad \omega\in\Omega_n.
\]
It is measurable, including when \(n=0\). For every family member and every sample,
\[
|\widehat\theta_{\mathrm{known},n}(\omega)-\mu_P(x_0)|=0.
\]
Since \(cr>0\), its miss event is empty. Consequently,
\[
\begin{aligned}
0
&\le
\inf_{\substack{\widehat\theta:\Omega_n\to\mathbb R\\
\widehat\theta\ \mathrm{measurable}}}
\sup_{P\in\mathfrak P}
P^{\otimes n}\{|\widehat\theta-\mu_P(x_0)|>cr\}
\\
&\le
\sup_{P\in\mathfrak P}
P^{\otimes n}
\{|\widehat\theta_{\mathrm{known},n}-\mu_P(x_0)|>cr\}
=0.
\end{aligned}
\]
This proves the asserted equality in the extended nonnegative reals.
% lean: CausalSmith.Stat.WeakOverlap.knownRegression_minimax_miss_zero

\item \emph{The causal-class pointwise lower bound.}
Apply the pointwise lower conclusion, supported by the two-valued potential-outcome subclass, of
\cref{obj:thm:adaptive-minimax} with the fixed parameters
\[
(B,L,C,c_f)=(B_*,L_*,C_{\mathrm{up}},1).
\]
Their required positivity and range conditions were established above. Distinguish the restricted pointwise risk by defining
\[
R_{\mathrm{pt}}^{\mathrm{bd}}(n,\gamma,x_0)
:=
\inf_{\widehat\theta\ \mathrm{sample\text{-}measurable}}
\sup_{\substack{P\in\mathcal P_\gamma\\
P\text{ admits a causal completion with}\\
Y(0),Y(1)\in\{0,B_*\}\ \text{a.s.}}}
\mathbb E_{P^{\otimes n}}
\left|\widehat\theta(\mathcal D_n)-\mu_{1,P}(x_0)\right|.
\]
For every \(\gamma>1\) and every \(x_0\in[0,1]^d\), the cited conclusion supplies \(c_\gamma>0\) such that
\[
c_\gamma r_{\gamma,n}
\le R_{\mathrm{pt}}^{\mathrm{bd}}(n,\gamma,x_0)
\le R_{\mathrm{pt}}(n,\gamma,x_0),
\qquad n\ge1.
\]
The second inequality follows because the restricted laws form a subset of the full causal class in \cref{obj:def:risks}: for each measurable estimator, restriction decreases the worst-case risk, and taking the infimum preserves this inequality. This gives the final assertion.
% lean: CausalSmith.Stat.WeakOverlap.adaptive_minimax_rate
\end{enumerate}
\end{proof}

\begin{proof}[Proof of \cref{obj:lem:dorn-rate-scope}]
\begin{enumerate}
\item Fix a finite \(\gamma_0>1\), the source family \(\mathfrak P\), its common constants, and its selected propensity representatives. The design, moment, variance, Gaussian, smoothness, and global-tail hypotheses give Dorn's Assumption 2, including the restrictions retained from Assumption 1. The family-level condition in \cref{obj:def:dorn-a3} gives Assumption 3 on the original cube.

Dorn's fixed-setup upper comparison \citep[Section 2, Assumptions 1--3, Theorem 1(ii)]{Dorn2026GlobalOverlap} therefore supplies an estimator \(\widehat\mu_n\), measurable as a function of the sample for every \(n\) and \(x\), with comparison scale
\[
r^*_{\mu,n}
=
n^{-\beta/(2\beta+d\gamma_0/(\gamma_0-1))}
\qquad(n\ge1),
\qquad r^*_{\mu,0}=0.
\]
Its admissibility also gives measurability, for every \(n\), \(P\in\mathfrak P\), and \(R\in\mathbb R\), of
\[
\left\{
\|\widehat\mu_n-\mu_P\|_{L^\infty(P)}
>
\max\{Rr^*_{\mu,n},0\}
\right\}.
\]
For every \(\varepsilon>0\), the cited comparison supplies a finite \(c(\varepsilon)>0\) such that
\[
\limsup_{n\to\infty}\sup_{P\in\mathfrak P}
P^{\otimes n}\!\left\{
\|\widehat\mu_n-\mu_P\|_{L^\infty(P)}
>
c(\varepsilon)r^*_{\mu,n}
\right\}
\le\varepsilon.
\]
The setup was fixed before selecting the estimator and threshold constant, so both may depend on that setup. This proves the upper comparison.
% lean: CausalSmith.Stat.WeakOverlap.DornSourceRateTranscription

\item Fix \(\gamma_0>1\). The hypotheses \(d\ge1\), \(\beta>0\), and \(L_0>0\) permit application of \cref{obj:lem:arbitrary-family-lower-obstruction}. Its known-regression conclusion gives, for every nonempty family with selected treated curves equal pointwise to \(\mu^\circ\), every evaluation point \(x_0\), every sample size \(n\), and every \(c,r>0\),
\[
\inf_{T\ \mathrm{sample\text{-}measurable}}
\sup_{P\in\mathfrak P}
P^{\otimes n}\!\left\{
|T(\mathcal D_n)-\mu_P(x_0)|>cr
\right\}
=0,
\]
with the minimax value taken in the extended nonnegative reals.

The singleton conclusion of the same result supplies \(q'>3\) and \(M',C'>0\) for the independent experiment displayed in the statement. It supplies the uniform design, selected propensity \(1/2\), moment and regression-variance bounds, zero treated and control curves with the required source smoothness, global-tail bound, Gaussian conditional outcomes, and the family-level condition of \cref{obj:def:dorn-a3}. The conditional Gaussian distributions have variance one, hence
\[
\operatorname{Var}_{P}(Y\mid X,A)=1=\sigma^2
\qquad P\text{-almost surely},
\qquad \sigma=1.
\]
Thus the conditional-variance restriction also holds for the asserted parameter tuple. Finally, the singleton conclusion at \(x_0=0\) gives
\[
\inf_{T\ \mathrm{sample\text{-}measurable}}
P^{\otimes n}\!\left\{|T(\mathcal D_n)|>cr\right\}
=0
\qquad(n\in\mathbb N,\ c,r>0).
\]
These are the two pointwise conclusions.
% lean: CausalSmith.Stat.WeakOverlap.arbitraryFamily_lower_obstruction
\end{enumerate}
\end{proof}
